\subsection{Local norms and level traces on the selected chain}
\label{subsec:selected-norms}

The reference projection fixes the row and its rank. Its height also depends on the local change of Schwartz vector and on the trace from one level to the other.

\begin{established}[Norm of the reference vector at a split place]\label{hyp:split-twist-norm}
At a split twist place \(v\) of a used tuple let \(q_v\) be the residue
cardinality, \(r\) the smaller block rank and \(b_v\) the larger rank.
Let \(L_0\) be the integral matrix lattice in the marked Schr\"odinger
model. With the unit reference vacuum, the local measure is
\[
 d\mu_v=dx/\operatorname{vol}_{dx}(L_0),\qquad
             \|\mathbf1_{L_0}\|_{L^2(d\mu_v)}^2=1 .
\]
\end{established}
\begin{proof}
The split places in the tuple construction avoid the dyadic
conductor, the marked auxiliary prime and the auxiliary ideal.
The chosen split projection identifies their residue field with the
one on the marked CM factor, and the hom-lattice ideal is a unit.
Thus the reference vacuum is the indicator of the integral lattice
in the split frames. Its squared norm for any Haar measure \(dx\) is
\(\int_{L_0}1\,dx=\operatorname{vol}_{dx}(L_0)\). Normalizing this norm
fixes the stated measure. All other level factors and all real data
agree by Lemma~\ref{hyp:joint-period-towers}.
\end{proof}

\begin{lemma}[Local norm of the split twist]\label{lem:split-twist-norm}
With the split model and measure of
Lemma~\ref{hyp:split-twist-norm}, the normalized local squared norm at a new split
twist place \(v\) is
\[
 q_v^{-r(b_v-r)}\prod_{j=1}^r(1-q_v^{-j}).
\]
\end{lemma}

\begin{proof}
In these coordinates the squared absolute value of the
twist vector is an indicator. All matrix entries are
integral. The first \(r\times r\) block is invertible
modulo \(v\), and the \(r\) entries of each of the
remaining \(b_v-r\) rows in those first columns are zero
modulo \(v\). There is a unit-modulus character of the
determinant on the invertible block. Counting independent
ordered columns in the \(r\)-dimensional residue vector
space shows that the invertible fraction among \(r\times r\)
matrices is
\[
             \prod_{j=1}^r(1-q_v^{-j}).
\]
The zero block contributes \(q_v^{-r(b_v-r)}\).
The other matrix coordinates contribute one in the
normalized measure. Thus the norm-square factor of the
twist at \(v\) is
\begin{equation}\label{eq:selected-twist-norm}
         q_v^{-r(b_v-r)}\prod_{j=1}^r(1-q_v^{-j}).
\end{equation}
\end{proof}

\begin{lemma}[Norm ratio and integral character trace]\label{lem:twist-norm-trace}
Let \(A\) denote the squared norm of a twisted step divided by that of its reference vector. With both sums over the new twist places,
\[
 \left|\log A+\sum_v r(b_v-r)\log q_v\right|\le2\sum_v(q_v-1)^{-1}.
\]
 At the prescribed finite level \(P\) the
full theta series is the character-weighted trace of the principal
series. At a coefficient place away from the tuple primes this trace
acts integrally with the chosen normalization.
\end{lemma}

\begin{proof}
The norm square of the global step is the product of
norm squares of the local factors. For the reference step
one uses the same factors at the common level primes and the
same real factors, and replaces each new twist factor by
the unit vacuum. For the common and twist prime sets
\(P_{\rm c},P_{\rm tw}\), the sum
\(e=\sum_{v\in P_{\rm c}\cup P_{\rm tw}} e_v\) of the chosen
split idempotents satisfies \(e\bar e=0\).
At a split real-field place the two primes above it give the
two complementary local idempotents. Indeed the local factor is a product of two fields and conjugation
exchanges the factors.

If the selected components there are marked
by \(a,b\in\{0,1\}\), then
\[
 e=(a,b),\qquad\bar e=(b,a),\qquad e\bar e=(ab,ab).
\]
Their product vanishes only when at most one component is selected. Thus
distinct primes in the union lie over distinct real-field places.
The ratio of squared norms is the product of
\eqref{eq:selected-twist-norm} over the new places. Taking
logarithms, the scalar correction beyond
\(-\sum_v r(b_v-r)\log q_v\) is bounded by
\[
 \sum_v\sum_{j=1}^r -\log(1-q_v^{-j})
  \le2\sum_v\sum_{j=1}^\infty q_v^{-j}
  =2\sum_v (q_v-1)^{-1}.
\]
The inequality uses \(q_v^{-j}\le1/2\) for every new place.
\label{cor:selected-norms-1}

\proofstep{The trace before division by the group order.}
The finite character trace uses the same level and tuple.
If \(P\) is the finite row level, \(\xi\) the character
on the lift, and \(a_g\) local integral representatives of
its elements, the prescribed full Schwartz vector is
\[
          \Phi=\sum_{g\in P}\xi(g)^{-1}
                         \omega(1,a_g)\Phi_{\rm pr}.
\]
Here \(\Phi_{\rm pr}\) is the principal vector. The character
of the full level on the lift is \(\psi_0^{-1}\). The
coefficient in the sum is its inverse. At a product of
places the trace is the product of the corresponding sums,
or equivalently a sum over tuples of local representatives
with the product character. This follows by expanding the
finite sums, since operators at distinct places commute.
The level lifts have the same residue entries as the
chosen finite group and the groups are identified by that
reduction at every selected place.

Write \(\Theta(\Phi)\) for the series with input \(\Phi\), and \(R(a_g)\) for right translation of its automorphic argument. Covariance at an individual level lift gives
\[
 \Theta\bigl(\omega(1,a_g)\Phi_{\rm pr}\bigr)=R(a_g)\Theta(\Phi_{\rm pr}).
\]
Linearity of the theta summation now gives the equality
of the full-vector series with the character-weighted
trace of the principal series, for each finite class and
real polynomial. There is no group-order denominator in
this equality. For the unnormalized character trace \(T_\xi\), translating the
summation variable gives
\[
 T_\xi=\sum_{g\in P}\xi(g)^{-1}R(g),\qquad
 T_\xi^2=|P|T_\xi .
\]
Here \(R(g)\) is right translation. Unitarity makes \(T_\xi\) self-adjoint, and group averaging makes the normalized operator an orthogonal projector:
\[
 T_\xi^*=T_\xi,\qquad P_\xi=T_\xi/|P|,\qquad P_\xi^2=P_\xi .
\] The integral
calculation uses the sum \(T_\xi\), including at primes dividing
\(|P|\).
\label{cor:selected-norms-2}

\proofstep{The scalar at the coefficient prime.}
For the chosen twist the comparison of principal vectors
is by a translation at the prescribed split places and a
scalar. The translated class is admitted because those
places are coprime to the excluded coefficient prime.
The factor that converts the principal vector and the
polynomial reading to the normalized row has valuation
zero there. In the chosen norm comparison a positive
power of the ratio of the two scalars is one. The separate
translation scalar has a positive power supported on the
tuple primes, which are coprime to the coefficient prime.

For either scalar \(z\) in the comparison take a positive power
whose prime support is contained in the tuple primes. At a coefficient
place \(w\) away from these primes,
\[
 n\operatorname{ord}_w(z)=\operatorname{ord}_w(z^n)=0.
\]
If the scalar ratio is read initially in \(\C\), a nonzero
coefficient in the faithful rational row comparison
realizes it over the enlarged coefficient field. The
vanishing case requires no scalar division. Hence, at the
finite places used in the twisted row lattice, the trace
calculation acts on the measured integral coefficients.
This is the trace input to the containment in
Section~\ref{subsec:twist-row-lattices}. It keeps the
reference normalization of the previous subsection.
\label{cor:selected-norms-3}
\end{proof}

\begin{proposition}[Norm cost of the ordinary and coupled chain steps]\label{prop:selected-step-cost}
Use the finite class, index and coupling bounds of
Proposition~\ref{prop:reference-class-data} and the single-kernel
logarithmic squared bounds \(k_j(v)\) at one embedding \(v\).
Put \(U_u=fm^2N_m^2\log(g_0g_1T_u)\). At each step the logarithmic squared-norm bound available for composition is
\[
 k^{\rm chain}_j(v)=k_j(v)+2(C_{\rm cl}+C_{\rm idx}+C_{\rm cos})U_u.
\]
This prescription includes both final lifts when \(m=2\).
\end{proposition}

\begin{proof}
The norm identity for the input Schwartz vectors is used
alongside an operator bound for the column chain. For the
latter, each matrix entry in the finite level model is a sum
over class and coupling labels. At an embedding the value
metrics are positive semidefinite. In the quotient by the null space, the contribution of \(h\)
images \(z_b\) to one value satisfies
\[
 \Big\|\sum_{b=1}^h z_b\Big\|^2\le
 \Big(\sum_b\|z_b\|\Big)^2\le h\sum_b\|z_b\|^2 .
\] Applying this at the output
class and then summing with the class weights bounds the
squared operator norm by the largest squared local
kernel bound times the two class counts and the two
index counts in the input and output. Replacing the
counts by their exponential bounds adds their logarithms
to the logarithmic kernel bound.

\proofstep{A single coupling coset for the last two steps.}
At the end the two steps share a coupling coset.
Fix an embedding and let \(i\) run over the input indices,
\(a\) over the classes at \(i\), \(b\) over the intermediate
classes, and \(c\) over the output classes. Write \(Y(b)\)
for the admissible coupling cosets at \(b\). With the
reciprocal stabilizer weights \(\alpha_{ia},\beta_b,\gamma_c\),
the terminal value at \(c\) is
\[
 \sum_{i,a}\alpha_{ia}\sum_b\beta_b
             \sum_{Y\in Y(b)} K^2_{i,b,Y,c} K^1_{i,a,b,Y} x_{ia}.
\]
Both kernels are read on the same coset \(Y\). Suppose the
squared bounds of their individual operator norms are
\(t_1,t_2\), and the class counts for the three spaces are
bounded by \(H_0,H,H_1\). Put \(I=\#\{i\}\) and
\(N_Y=\max_b |Y(b)|\), taking \(N_Y=0\) for an empty
intermediate set. The weights are at most one.

Taking
the norm in the output fibre at \(c\), the triangle
inequality and the kernel bounds give the upper bound
\(\sqrt{t_1t_2}\,H N_Y\sum_{i,a}\alpha_{ia}\|x_{ia}\|\).
Cauchy--Schwarz bounds the square of the last sum by
\(I H_0\sum_{i,a}\alpha_{ia}\|x_{ia}\|^2\).
Summing over \(c\) with weights \(\gamma_c\) therefore bounds
the squared norm of the coupled operator by
\begin{equation}\label{eq:selected-coupled-norm}
       I H_0 H^2 H_1 N_Y^2 t_1t_2.
\end{equation}
Each calculation holds for the semidefinite metrics by passing
to their null quotients. The factor \(N_Y^2\) accounts for the
sum over the single shared coset in the terminal value.

\proofstep{Insertion of the class counts.}
On a selected sequence with prime \(T_u\), let \(N_m\) be the
last unitary rank, \(f\) the number of real-field places
and \(g_0,g_1\) the type primes. Use the common unit
\[
        U_u=fm^2N_m^2\log(g_0g_1T_u).
\]
Proposition~\ref{prop:reference-class-data} gives the class,
index and coupling counts on the chosen frames, with
upper bounds \(e^{C_{\rm cl}U_u},e^{C_{\rm idx}U_u}\) and
\(e^{C_{\rm cos}U_u}\), respectively.

Write \(k_j(v)\) for
the logarithmic bound for a single kernel of lift \(j\)
at the embedding \(v\). The ordinary step then has
logarithmic squared bound
\(k_j(v)+2(C_{\rm cl}+C_{\rm idx})U_u\).
The coupled bound \eqref{eq:selected-coupled-norm} has
logarithm at most
\[
 k_{m-1}(v)+k_m(v)
        +(4C_{\rm cl}+C_{\rm idx}+2C_{\rm cos})U_u .
\]
Both estimates are covered by enlarging each single-kernel
bound by \(2(C_{\rm cl}+C_{\rm idx}+C_{\rm cos})U_u\).
The admissible terminal data supply the two kernel bounds
with these indices through the first coupled entry and the
last padded entry. Composition adds these logarithmic bounds at
the same embedding, including both end lifts even when \(m=2\).
This estimates the operator in the reference-chain coordinates
of \eqref{eq:selected-pairing}.
\label{cor:selected-norms-4}
\end{proof}

\subsubsection{Counting the classes of the reference chain}
\label{subsec:selected-class-counts}
\hypertarget{source-selected-class-counts}{}

The reference norm is a sum over definite classes. Bounding the number of its terms requires a class count for the ranks occurring along this chain.

Fix the reference tuple, the prime assigned to the sequence \(T=T_u\), and the two
type primes \(g_0,g_1\) for the chosen CM extension \(E/F\). Write
\(O=\mathcal O_E\), \(d=[E:\mathbb Q]=2f\), and
\(r=|D_E|^{1/d}\). Let \(n\) be a rank occurring in the count,
\(1\le n\le N_m\). We first bound the number of rational unitary
orbits of self-dual \(O\)-lattices in the standard positive Hermitian
space \((E^n,h)\).
\label{cor:selected-class-counts-1}

\begin{proposition}[Size of the reference class data]\label{prop:reference-class-data}
For the selected reference tuple on \(E/F\), with type primes
\(g_0,g_1\) and assigned prime \(T_u\),
use the logarithmic size unit \(U_u=fm^2N_m^2\log(g_0g_1T_u)\).
Write \(h,i,s\) for the class, matrix-index and intermediate-coset counts. Then
\[
 h\le e^{55U_u},\qquad i\le e^{U_u},\qquad s\le e^{2U_u}.
\]
 With the
metrics of the selected chain, their logarithmic squared-norm
costs are at most \(112U_u\) in an ordinary step and
\(225U_u\) in the two coupled end steps together.
\end{proposition}

\begin{proof}
\proofstep{Find a short basis by transference against the trace dual.}
Give a lattice \(M\) the real inner product
\(\operatorname{Tr}_{E/\mathbb Q}h(x,y)\), with the convention that
the trace sums over all embeddings. Its integral trace dual is
\(\mathfrak D_E^{-1}M\). Indeed the trace pairing is nondegenerate
on the rational space. Testing a candidate dual vector against all
\(O\)-multiples of every vector of \(M\) puts its Hermitian
pairings in the codifferent, and localization of Hermitian
self-duality gives the displayed module.

If \(x\ne0\) belongs to
this dual, \(h(x,x)\) is totally positive and belongs to
\(\mathfrak D_E^{-2}\), since the different is conjugation stable.
Thus \(|N_{E/\mathbb Q}h(x,x)|\ge |D_E|^{-2}\). Since \(d=[E:\Q]\) and \(r=|D_E|^{1/d}\), applying the arithmetic
and geometric mean inequality to the conjugates gives
\[
 \|x\|^2=\operatorname{Tr}_{E/\Q}h(x,x)
     \ge d\,|N_{E/\Q}h(x,x)|^{1/d}\ge d/r^2.
\]
The Euclidean transference theorem in dimension \(dn\) bounds the
last successive minimum of \(M\) by \(n\sqrt d\,r\). These minima
are attained by discreteness. Choose an \(E\)-basis \(v_1,\ldots,v_n\)
from a rationally spanning family attaining this bound.
\cite{banaszczyk}

\proofstep{Count Gram matrices and the lattices between two integral bounds.}
Put \(B=4f^3 n^4 r^2\). The Gram matrix \(A=(h(v_i,v_j))\) has
entries in \(O\). Cauchy--Schwarz at each embedding bounds their
absolute values by \(dn^2r^2\le B\). The unit ball for the maximum
of the archimedean absolute values has real dimension \(d\).
Distinct elements of \(O\) have distance at least one in this norm
by the product formula. Packing translates of its open half ball
bounds the number of choices per entry by \((1+2B)^d\).
In the \(v_i\) basis, with the form written \(\bar x^t A y\),
\[
           O^n\subseteq M\subseteq A^{-1}O^n .
\]
The order of the quotient of the outer modules is
\(|N_{E/\mathbb Q}\det A|\le(nB)^{dn}\) by Hadamard.

Each submodule of this quotient is generated by at most \(n\)
elements. Locally its inverse image in \(A^{-1}O_P^n\) is a
lattice free of rank \(n\) over the DVR. Pad the lists of
generators of its images to length \(n\), and combine them at the
finitely many primes of the quotient using the Chinese remainder
theorem. If \(H=A^{-1}O^n/O^n\), the generator lists give a surjection
\[
 H^n\longrightarrow\{U:O^n\subseteq U\subseteq A^{-1}O^n\},
 \qquad (x_1,\ldots,x_n)\longmapsto
  \text{the inverse image of }\sum_i O x_i .
\]
Hence the count is at most \(|H|^n\le(nB)^{dn^2}\).
The pair consisting of \(A\) and this coordinate module determines
one rational unitary class, by mapping an ordered basis to another
basis with the same Gram matrix.

Set \(P=d^2 n^4 r^2\), so \(B=fP\). As \(B=fP\ge1\) and \(f^2n\le P\), the matrix count times the
module count is bounded by
\[
 (1+2B)^{dn^2}(nB)^{dn^2}
    \le (3f^2nP^2)^{dn^2}\le(3P)^{3dn^2}.
\] This
argument allows projective \(O\)-lattices of every Steinitz class.

At a selected place \(v\) over \(T\) a decoration specifies a
residue functional. At a selected place over a type prime it
specifies \(n\) such functionals, the coordinates of its frame.
On a fixed lattice of rank \(n\), each functional is determined by
its values on lifts of a residue basis. The localized lattice
modulo its maximal ideal has dimension \(n\). One can take the
lifts in the global lattice since denominators outside the place
are invertible on the residue field. The rules for a decoration
determine it by restriction to the localized lattices, including
when its definition represents the functions on the ambient space.

Writing \(q_v\) for the residue cardinality, we have at most
\[
        \prod_{v\in S_T}q_v^n\prod_{v\in S_G}q_v^{n^2}
          \le (g_0g_1T)^{fn^2}
\]
choices. The selected half above each of the three split primes has
\(f\) places of residue degree one. Here \(S_T\) can be a subset
of the word half and \(S_G\) is the union of the two type halves.
Transport each decorated point to one of the underlying lattice
representatives by a rational unitary transformation, leaving the
external frame-change group element equal to the identity. For a fixed underlying lattice each class costs at most the decoration count just obtained. Summing gives
\[
 \#\{\text{principal decorated classes}\}\le \#\{\text{underlying lattice classes}\}\,\max_L\#\{\text{decorations on }L\}.
\] Stabilizers can only reduce this number.

We explain how the same bound applies to each of the three class
spaces in the period chain. Suppose a point space injects into the
decorated points and every action of \((g,1)\), for \(g\) rational
unitary, on an image point lifts to an allowed move. Choose one
representative per orbit in the point space. If two chosen images
belong to one principal class, lift the relating move and apply
injectivity. Concretely, for the injection \(i\) on points and its allowed lifted
move \(\tilde g\), one has
\[
 i(y)=(g,1)i(x)=i(\tilde g x)\quad\Longrightarrow\quad y=\tilde g x .
\]
Thus taking principal classes does not identify distinct point-space
orbits.

In a class space of
a column of rank \(N_j\), forgetting the supplementary twist
labels is injective because the base tuple has empty twist
support. The move lifts as \((g,1,1)\) at every matrix index.
For the evaluation rows take rank \(N_m-1\), empty word decoration,
and the identity on decorated points. The move is \((g,1)\).
For the top class count take rank \(N_m\), the full word decoration
and the identity, with \((g,1)\) in the permitted product subgroup.
The two type halves exhaust the row places at the base tuple. Thus
all three counts use precisely the prime product \(G=g_0g_1T\).

The selected size conditions give the following simultaneous range:
\[
 L\ge1,\qquad\log r\le L,\qquad fN_m\le e^L\le g_0 .
\] In this comparison it is
enough to use these four inequalities. Since \(12\le e^{3L}\),
\[
         3P=12(fn)^2 n^2 r^2\le e^{9L}\le G^9 .
\]
The resulting bound on each of the three class counts is
\[
 G^{(54+1)fn^2}\le\exp(55U_u),\qquad U_u=fm^2N_m^2\log G .
\] This is uniform in the base row,
column, index and word. \label{cor:selected-class-counts-2}

\proofstep{Read each index on the residue bases of the padded lattices.}
For the matrix indices only stages \(j+2\le m\) occur. The
initial index is unique. For a positive stage record the residual
matrix of its representative at each type place on the padded
lattice of rank \(n'=N_j+c_j\le N_m\). In coordinates with lattice
exponents \(e_a\) (zero on the initial block and successive
positive integers on the padding), its entry is the reduction of
\(g_i^{e_b-e_a}h_{ab}\) at a place over \(g_i\).

In these coordinates the residual matrix is the reduction of
\(D^{-1}hD\), with \(D=\operatorname{diag}(g_i^{e_a})\). Equal
residual matrices for \(h_1,h_2\) imply
\[
 D^{-1}h_1^{-1}h_2D
     =(D^{-1}h_1D)^{-1}(D^{-1}h_2D)\equiv I\pmod{g_i}.
\] The type prime is unramified of residue degree one.
Every defining flag lattice lies between the padded lattice and
its \(g_i\)-multiple, and the prescribed residue vectors lie in the
padded lattice. Hence the relative action preserves each flag
lattice and fixes the vector classes in its prescribed quotients.
This is the right stabilizer in the local double coset. Equality
of these double cosets at all the selected places identifies the
index by its defining injectivity. At the base tuple the resulting count of indices satisfies
\[
 \#\{\text{indices}\}\le(g_0g_1)^{f(n')^2}\le e^{U_u}.
\]

\proofstep{Bound the cosets by the single modulus at the two type primes.}
The intermediate coupling set consists of cosets of a module
\(A_0\) represented by elements of a module \(A_1\subset E^n\).
The defining local conditions give \(g_0g_1 A_1\subseteq A_0\):
at a place over \(g_i\) the one relaxed lattice exponent is
restored by multiplication by \(g_i\), and there are no twist
places. To count the quotient, let \(L\subset A_1\) be generated by
finitely many representatives. The inclusion just proved gives
\[
 L/(g_0g_1)L\twoheadrightarrow L/(L\cap A_0),\qquad
 |L/(L\cap A_0)|\le (g_0g_1)^{\operatorname{rank}_{\mathbb Z}L}
                 \le (g_0g_1)^{2fn}.
\]
Every finite set of cosets is counted in such an image. This proves the bound even for the semilocal modules.

The coset bound applies to every actual intermediate lattice as well as on the
standard one. We may take \(C_{\rm cos}=2\),
\(C_{\rm idx}=1\), and \(C_{\rm cl}=55\). In particular the
ordinary and coupled squared-norm losses above are at most
\(112 U_u\) and \(225 U_u\), respectively. Increasing the
single-kernel constant by \(58\) pays \(116U_u\) per lift.
Both end lifts, at the same word and embedding, are charged
in the coupled case. \label{cor:selected-class-counts-3}
\end{proof}


\subsubsection{The archimedean jet of a single kernel}
\label{subsec:selected-kernel-jets}

For a single term in the class sum one must bound a derivative of the archimedean kernel, including both jets when the last two lifts are coupled.

\paragraph{The padded polynomial bound.}

The first estimate bounds the archimedean polynomial after extending it to the larger of the two ranks in a lift.

\begin{established}[Determinant shifts in the padded jet model]\label{hyp:padded-Gaussian-jets}
For integers \(1\le n\le n'\), let \(\nu_\tau\) be a partition of
\(q_\tau\) of length at most \(n\) at each real label \(\tau\).
On the source divide out the common last entry, and on the target
use the partition padded with zeros to length \(n'\), up to its
determinant character. With the highest-vector metrics, restoring
the source determinant factor by multiplication with
\(\det^{\nu_{\tau n}}\) gives an isometry \(T\). In particular
\[
               T^*T=I,\qquad \deg(Tv)=|\nu_\tau|=q_\tau .
\]
\end{established}
\begin{proof}
In the flag-minor model the determinant has weight \((1,\ldots,1)\)
and has absolute value one on the unitary group. Thus \(T^*T\)
commutes with that group. Schur's lemma makes it scalar, and the
normalized source highest polynomial is sent to the normalized
target highest polynomial. The scalar is their squared norm ratio,
which is one. If \(n=n'\), padding inserts no zeros. The common
last entry is removed by the determinant twist in this case as well.
\end{proof}

\begin{lemma}[Norm of a padded jet contraction]\label{lem:padded-jet-estimate}
Use the polynomial spaces of Lemma~\ref{hyp:padded-Gaussian-jets}.
For \(Z_\tau\in\mathbb C^{n'\times n}\) substitute
\(Y_\tau=Z_\tau^t X_\tau E_n\), where \(E_n:\mathbb C^n\hookrightarrow\mathbb C^{n'}\)
includes the first coordinates and \(X_\tau\) is the square target
variable. Let \(G\) be entire, with \(S\ge0\), \(r_\tau,\Omega_\tau>0\) and
\[
 |G(Z)|\le S\exp\!\left(\tfrac12\sum_\tau r_\tau^2\|Z_\tau\|^2\right),
 \qquad |r_\tau d_\tau|^2=\Omega_\tau^{-1}.
\]
Here the matrix norm is Euclidean on the entries.
The jet replaces \(Z^\alpha\) by
\((\prod_\tau d_\tau^{|\alpha_\tau|})\partial^\alpha G(0)\).
The squared operator norm of this contraction,
in the normalized highest-vector metrics, is at most
\begin{equation}\label{eq:selected-padded-constant}
 S^2\prod_\tau\left[
    \Omega_\tau^{-q_\tau} e(2q_\tau+2)^{2nn'}
    \exp\!\left(\sum_i\varphi(\nu_{\tau i})\right)
 \right].
\end{equation}
This bound includes labels of degree zero.
\end{lemma}

\begin{proof}
Write \(W_\nu(\C^d)\) for the Schur module with its unitary
metric, and \(v_{\nu,d}\) for its initial flag vector of norm
one. The realization \(u\mapsto u^t W_\nu(Y)v_{\nu,d}\)
is the flag-minor model, with the transpose taken in bases
compatible with the real structures. In fact the highest
matrix coefficient is the product of initial flag minors,
and the row action generates the model. The Schur maps
commute with transpose in these bases. Functoriality and compatibility with transpose give
\[
 u^tW_\nu(Z^t X E_n)v_{\nu,n}
 =(W_\nu(Z)u)^t W_\nu(X) W_\nu(E_n)v_{\nu,n}
 =(W_\nu(Z)u)^t W_\nu(X)v_{\nu,n'} .
\] The identity uses
\(W_\nu(E_n)v_{\nu,n}=v_{\nu,n'}\), with compatible highest
vectors. Thus the coefficient contraction lands in the
target highest-weight model. This explains the order
\(Z^tX E_n\) in the substitution and the absence of a
projection cost.

\proofstep{Gaussian contraction of the two representations.}
Use the Gaussian polynomial norm
\(\|p\|_{\rm F}^2=\sum_\alpha\alpha!|p_\alpha|^2\).
For rectangular \(Z\), a matrix coefficient
\(p(Z)=w^*W_\nu(Z)u\) has
\begin{equation}\label{eq:selected-fock-matrix}
 \|p\|_{\rm F}^2=H(\nu)\|w\|^2\|u\|^2,\qquad
 H(\nu)=\prod_{(i,j)\in\nu} h_{ij}.
\end{equation}
To see this, realize both Schur modules in tensor powers of
degree \(q=|\nu|\) by Schur--Weyl duality, using the same
unit vector \(s\) in the Specht factor of shape \(\nu\).
The operator \(Z^{\otimes q}\) acts as \(W_\nu(Z)\) on
the Schur factor and as the identity on that Specht factor.
Integrate a coefficient of \(Z^{\otimes q}\) times its
conjugate against independent standard complex Gaussians.

The moment formula pairs positions by all \(\pi\in S_q\),
giving the factor
\(\sum_\pi |\langle s,\pi s\rangle|^2=q!/\dim[\nu]\).
Indeed each pairing identifies the entry in position \(k\)
with the conjugate entry in position \(\pi(k)\), both in
its row and in its column. Indeed the average of the conjugates of the projector onto \(s\)
is scalar on the irreducible Specht module. Its trace is one, so
\[
 \frac{1}{q!}\sum_{\pi\in S_q}\pi P_s\pi^{-1}
     =\frac{1}{\dim[\nu]}I .
\]
Pairing both sides with \(s\) gives the asserted sum. The hook-length
formula gives \(q!/\dim[\nu]=H(\nu)\). This argument also
shows that the Fock normalizer on the square model for the
partition itself is \(H(\nu)\).
\cite{frame-robinson-thrall}

\proofstep{Bounding the homogeneous coefficient.}
For each label rescale to \(z_\tau=r_\tau Z_\tau\) and put
\(g(z)=G((z_\tau/r_\tau)_\tau)\). With \(D=nn'\), take
the torus Fourier projection of \(g\) onto multidegree
\((q_\tau)_\tau\), denoted by \(g_{\mathbf q}\). On the
product of spheres of radius \(\sqrt{q_\tau}\) its
absolute value is at most \(S e^{\sum_\tau q_\tau/2}\).
Use radius zero for a zero degree. Gaussian polar
integration of this homogeneous component yields
\[
 \|g_{\mathbf q}\|_{\rm F}^2
 \le S^2\prod_{\tau:q_\tau>0}
 \frac{e^{q_\tau}(q_\tau+D-1)!}
      {q_\tau^{q_\tau}(D-1)!}.
\]
The radial moment at one label is
\((q+D-1)!/(D-1)!\) because the squared radius of a standard
complex Gaussian vector in \(D\) dimensions has gamma
density with shape \(D\).

For \(q=0\) the factor is one.
For \(q>0\) the factor is at most \(e(q+1)^D\).
Indeed monotonicity of the logarithm gives
\[
 \log q!\le\log q+\int_1^q\log t\,dt
            =\log q+q\log q-q+1 ,
\]
so \(e^q q!/q^q\le eq\), and
\(\binom{q+D-1}{D-1}\le(q+1)^{D-1}\) by recording the first
\(D-1\) entries of a weak composition of \(q\).

\proofstep{Tensor products over the labels.}
For general vectors in the tensor product over labels the
Fock identity \eqref{eq:selected-fock-matrix} tensorizes.
For example, polarizing its proof gives the product of the
input and output Hermitian pairings for two matrix
coefficients. Applying Fock Cauchy--Schwarz to the bilinear
contraction with \(g_{\mathbf q}\) therefore gives the
squared bound \(\|g_{\mathbf q}\|_{\rm F}^2
\prod_\tau H(\nu_\tau)\Omega_\tau^{-q_\tau}\).

To bound \(H(\nu)\), its row product formula gives
\[
 H(\nu)=\frac{\prod_i(\nu_i+n-i)!}
               {\prod_{i<j}(\nu_i-\nu_j+j-i)}
 \le \prod_i \nu_i!\,(q+1)^{n(n-1)/2}.
\]
Here each denominator factor is at least \(j-i\), and
\(\binom{\nu_i+n-i}{n-i}\le(q+1)^{n-i}\).
Concavity of \(\log\) gives
\(k!\le e\sqrt{k}\,\exp(\varphi(k))\) for \(k\ge1\)
by the trapezoid lower bound on its integral from \(1\)
to \(k\). Thus
\(\nu_i!\le(2q+2)\exp(\varphi(\nu_i))\), including
\(\nu_i=0\). Since \(n(n+1)/2\le nn'\), multiplying
these bounds proves \eqref{eq:selected-padded-constant}.
\end{proof}

\paragraph{The finite datum and the two levels.}

The real estimate applies termwise. To use it in the class sum we must specify the finite levels on the two sides of the same entry.

Use the untwisted reference tuple, and fix a word \(u\)
and the two classes of one matrix entry. At the coupled
step also fix their common coset. Put \(g=g_0g_1\),
\(c=T_u g\), and write \(c_B=cc'\). For the integral level exponents \(A,B\) the reference scalars are
\[
 c'=g^{fB-1},\qquad c_A=g^{fA},\qquad fB\ge1 .
\]
Let \(H\) be the semilocal integral Hom module on the two
class lattices at the selected places. In the actual
finite datum \(\Xi\), the factor at the assigned prime is the
indicator of the integral Hom condition there. For a type place \(v\), write \(H_v\) for its integral Hom lattice, \(\mathfrak p_v\) for the maximal ideal and \(f_v\) for the local factor. The two lattice frames give 
\[
 \operatorname{supp} f_v\subseteq H_v,\qquad f_v(x+y)=f_v(x)\quad(y\in\mathfrak p_vH_v).
\]

In ordinary
steps this follows by unfolding the extra columns with
the level character. For a place of residue size \(Q\),
the character has conductor one. If \(u_i,v_i\) are
the ranks including the extra rows, and \(s_i\) counts
extra source columns, every free coordinate in a
contributing matrix has valuation at least \(-s_i\).
At reading depth \(B\) it ranges over at most \(Q^{B+s_i}\)
residue choices. For each free entry the contribution from counting residues and
integrating their cells is bounded by
\[
 Q^{B+s_i}\operatorname{vol}(\varpi^B O)=Q^{B+s_i}Q^{-B}=Q^{s_i},
\]
where \(\varpi\) is a uniformizer and \(O\) has measure one.
The character on the shifted diagonal has argument integral after
multiplication by \(Q\). Hence its value has absolute value one at
every embedding. Multiplying the coordinate estimates and the inverse volume of
the output unipotent gives
\[
 (Q^{B+s_i}Q^{-B})^{u_i v_i}\,Q^{v_i^2}
     =Q^{s_i u_i v_i+v_i^2}\le Q^{2m^2N_m^2}.
\]

Here \(u_i,v_i\le N_m\), \(s_i+1\le m\), and the last factor
on the left is the subgroup index in the padded integral matrix space. In the first coupled entry the
prescribed integral column is inserted into the moved
matrix test. In the last entry the top matrix test is
multiplied by the normalized finite Fourier test on
the last row. Both have residue invariance and bound
one, the latter by the triangle inequality in the
normalized residue average. In both cases the coset
is integral in its decorated frame. The residue sizes
at the type places multiply to \(g^f\). Therefore \(\Xi\)
vanishes outside \(H\), is constant on classes modulo
\(cH\), and, for an absolute \(C\), obeys the uniform bound
\[
      |\rho(\Xi(x))|\le\beta,\qquad \beta=g^{C f m^2 N_m^2}.
\]
\label{cor:selected-kernel-jets-1}

\paragraph{The coefficient field of the two jets.}

A coupled entry involves two jets at once. Their Galois conjugates can be averaged together only after both jets have been placed in one number field.

\begin{lemma}[One field for the two jets in the coupled entry]\label{lem:coupled-jets-field}
At the selected ranks and reference modulus \(M_0=c_Ac_B\),
the chosen coefficient field contains the jets at every permitted
class pair and indicator at each step, including both steps of a
coupled entry. Writing \(K\) for the chosen field and \(a,b\) for any two coupled jet entries, the field of definition of both steps already satisfies
\[
 a,b\in K,\qquad ab\in K.
\]
Their products and class contractions therefore precede the choice of archimedean embedding.
\end{lemma}

\begin{proof}
For the selected ranks \(a,b\), bad-place set \(S_b\), and
positive modulus \(M_0=c_Ac_B\), consider the normalized
indicator indexed by self-dual lattices \(L,L'\),
\(c\mid M_0\), and \(x\in M_0^{-1}H_{S_b}(L,L')\).
Here \(H_{S_b}\) is the semilocal Hom lattice in which
the lift's indicator is defined. For fixed \(L,L',c\),
the indicator depends on \(x\) modulo \(cH_{S_b}\).
The quotient \(M_0^{-1}H_{S_b}/cH_{S_b}\) is finite.

For \(H=H_{S_b}(L,L')\), multiplication by \(M_0\) gives
\[
 M_0^{-1}H/cH\xrightarrow{\ \sim\ }H/(cM_0 H),\qquad
     x+cH\longmapsto M_0 x+cM_0 H.
\]
The local quotient vanishes where \(cM_0\) is a unit.
Use one representative of each of the finitely many
self-dual lattice classes at both ranks. For each pair
and each divisor of \(M_0\), take the normalized Taylor
coefficients of the resulting finite list of indicators
through the total degree of the stage. These coefficients
are algebraic by the CM Taylor comparison on the even-rank
standard lattices and their translates. The hypotheses
are the fixed nonzero period, the compatible ideal character,
and split bad places where the period ideal is a unit.
These hold for the principal curve ideal and the selected
reference sequence.

\proofstep{Changes of representatives.}
Let \(g,h\) identify the source and target lattice representatives with the actual lattices. Composition with these identifications acts on Hom variables by
\[
 x\longmapsto h\circ x\circ g^{-1}.
\] This gives a bijection of the finite semilocal
quotients, since it identifies both \(M_0^{-1}H_{S_b}\)
and \(cH_{S_b}\). Equivariance substitutes the resulting
linear change \(L\) in the labelled Taylor variables.
For a homogeneous term of total degree \(d\), write the
normalized polynomial at the representative as
\(\sum_{|\beta|=d} a_\beta Z^\beta\).

For a coefficient with index \(\alpha\) the transported
reading is
\[
      d_\alpha\sum_{|\beta|=d}a_\beta[Z^\alpha](LZ)^\beta .
\]
Here \(d_\alpha\) is the product of the differential rescaling
in degree \(\alpha\) with the algebraic character ratio.
The latter is a product of determinants to integral powers.
The entries of \(L\) and those determinants use the labelled
images of the auxiliary field. The differential rescalings
use powers of the prescribed generators. Let \(F\) contain the representative coefficients, the entries
of \(L\), the determinant factors and the differential generators.
The multinomial theorem gives
\[
 [Z^\alpha](LZ)^\beta\in \mathbb Z[L_{ij}]\subset F
 \quad\Longrightarrow\quad
 d_\alpha\sum_{|\beta|=d}a_\beta[Z^\alpha](LZ)^\beta\in F.
\] In particular, the coefficients of multidegree
\((|\nu_\tau|)_\tau\), of total degree \(\sum_\tau|\nu_\tau|\),
belong to \(F\).

\proofstep{One finite set for both steps.}
The coefficient-field choice in the selected sequence family
uses this finite generation before the arithmetic row
and column are realized. At fixed predata there are
finitely many words and stages, and finitely many
bad places from the conductors of the character family
and the selected levels. Include the finite generators
for all the permitted subsets of these places.

For the
canonical twisted choice, if \(Q\) is the product of the
row prime norms, its two half moduli are
\(Q^{m+1}\) and \(T_uQ^{m+1}\). Thus the modulus of the stage list is precisely the product
\[
 Q^{m+1}(T_uQ^{m+1})=T_u Q^{2m+2}.
\] The family permits
requiring any such finite collection of algebraic
constants in the embedded coefficient field. Its
construction imposes this requirement jointly with the
column conditions. Thus the chosen field contains the
jets for the actual pair and indicator of each step.
In the coupled entry both consecutive kernels have
their own stage data in this same list. Their products
and the finite class contractions are therefore evaluated
in this field before taking any embedding.
\label{cor:selected-kernel-jets-2}
\end{proof}

\begin{proposition}[Coefficient mass in the torsion presentation]\label{prop:torsion-jet-mass}
At a reference level \(c\) of the selected kernel, use the
torsion presentation of the normalized entry and its finite-data
bound \(\beta\). If \(a_\nu\) are its torsion-presentation coefficients, then for each embedding \(\sigma\)
\[
 \sum_\nu|\sigma(a_\nu)|\le\beta\,c^{2[E:\Q]nn'}.
\]
 This holds also
for a vanishing entry.
\end{proposition}

\begin{proof}
In that field the entry is a jet sum over
\((c_A^{-1}H)/(c_B H)\).
Let \(F_\ell(x)\) be the normalized theta function
attached to the indicator of \(x+\ell H\), using
the class-pair vacuum. For the reference levels, the
odd integer \(c_Ac_B\) is supported on the specified
split half places, where the period ideal is trivial.
To compare the readings in \eqref{eq:unserved-reader}, write
\(A,B\) for the half lattices at \(cc'\), and \(A_1,B_1\)
for those at \(c\). The complementary exponents give
\[
 A\cap B_1\subseteq\Lambda,\qquad
 B_1=(A\cap B_1)+(B\cap B_1),\qquad
 \Lambda\cap A\subseteq A_1 .
\]

In the middle equality the expanding half of each conjugate pair
supplies the local lattice. Put \(D_f=\Lambda\cap A\cap B\)
and \(D_c=\Lambda\cap A_1\cap B_1\). The complementary exponents
also give \(A\cap B\subseteq B_1\), hence \(D_f\subseteq D_c\).
The inclusion \(\Lambda\cap A\subseteq\Lambda\cap A_1\) induces the finite map
\[
 (\Lambda\cap A)/D_f\longrightarrow (\Lambda\cap A_1)/D_c,
 \qquad t+D_f\longmapsto t+D_c .
\]
This map is onto. Passing from \(A_1\) to \(A\) imposes stronger
divisibility only where \(B_1\) expands; there the quotient allows
zero as representative. At the other primes keep the prescribed
residue. The Chinese remainder theorem gives one lift with all
these local residues.

\begin{construction}
\[
 \begin{array}{ccc}
 H/(cc'H)&\longrightarrow&H/cH\\
 x+y&\longmapsto&x \pmod{cH}\\[2pt]
 \{F_{cc'}(x+y):y\in cH/(cc'H)\}&\xrightarrow{\ \mathrm{sum}\ }&F_c(x)
 \end{array}
\]
\caption{Coarsening the Fourier readings in \eqref{eq:selected-kernel-regroup}. The lower arrow sums the fine readings in a single fibre of the upper quotient map.}\label{fig:fine-coarse-fourier}
\end{construction}

\proofstep{Regrouping the finite Fourier series.}
First read the indicator of \(B_1\) in the fine Fourier sum
with no outer translation. The indicator restricts the input
sum to the subgroup \((A\cap B_1)/(A\cap B)\).
For \(y\in B\) its alternating character \(\chi_y\) is already
trivial on \(B\cap B_1\). The sum decomposition of \(B_1\) above
and the self-duality of this lattice therefore give, successively,
\[
 \chi_y|_{A\cap B_1}=1
 \quad\Longleftrightarrow\quad \chi_y|_{B_1}=1
 \quad\Longleftrightarrow\quad y\in B_1 .
\] Orthogonality now leaves exactly the
coarse output representatives in \(B_1/(B_1\cap\Lambda)\).
The two finite groups in the surjection have orders \(h_f,h_c\),
the respective Fourier denominators in \eqref{eq:unserved-reader};
each fibre therefore has size \(h_f/h_c\).
Restrict \(\chi_y\) to
\(B_0=(A\cap B_1)/(A\cap B)\). Then
orthogonality gives
\[
 \frac{1}{|B_0|}\sum_{t\in B_0}\chi_y(t)=\mathbf1_{\{y\in B_1\}} .
\]

At a surviving frequency, counting the fibres of the preceding
surjection gives the coarse normalization \((h_f/h_c)/h_f=1/h_c\).
The semicharacter times the Weyl
translate of the vacuum depends only on the output class,
and the period and volume factors are common.

Expanding the fine reading by its input cosets is thus the
coarse reading at the origin. For a rational lift \(s\) of
the centre, reindex the fine input classes by translation
by \(s\) and act on both sides by \(\alpha(-s)U_{-s}\).
The Weyl composition law and the semicharacter law on \(A\)
give the translated fine indicators with
their same normalization. In semilocal coordinates this is
the partition of the coset modulo \(cH\), which proves
\[
 F_c(x)=\sum_{y\in cH/(cc'H)}F_{cc'}(x+y).
\]
Using the support and constancy of \(\Xi\) gives the
identity of entire sums
\begin{equation}\label{eq:selected-kernel-regroup}
 \sum_{x\in c_A^{-1}H/(cc'H)}\Xi(x) F_{cc'}(x)
        =\sum_{x\in H/cH}\Xi(x) F_c(x).
\end{equation}

In particular the larger denominator bound in the
coefficient field and the refinement factor do not
enter the archimedean count. The projective local
Hom modules have rank \(nn'\) over the local integer
rings, so
\(\#(H/cH)\le c^{[E:\Q]nn'}\). Writing \(q_w\) for the residue size, the count is the partial
product in the norm of \((c)\):
\[
 \prod_{w\ {\rm selected}}q_w^{nn' v_w(c)}
   \le\prod_{w\mid c}q_w^{nn' v_w(c)}
   = |N_{E/\Q}(c)|^{nn'} .
\] The missing factors, when present,
are positive integral powers in the full norm of \(c\).

\proofstep{Bounding the number and sizes of the coefficients.}
At level \(c\), write \(s\) for the rational lift of \(x\)
in the half lattice and use \(S_A,S_B\) as in the
finite Fourier reading. The expansion of \(F_c(x)\)
uses translates at \(y-s\), \(y\in S_B\). With \(v\) the
volume in the indicator transform, the two factors
\(v^{-1/2}\) cancel the \(v\) in its Fourier measure.
The phase multiplying the average at \(y\) is
\(\alpha(s)\alpha(y)e^{\pi i E_H(s,y)}\), since the
indicator translation is \(U_{-s}\). Write \(B_{\rm lat}\)
for the contracting half lattice. The average is
\[
 \frac1{|S_A|}\sum_{t\in S_A}{\bf1}_{B_{\rm lat}}(t)
                                e^{2\pi i E_H(t,y)}.
\]
All its phases, including the prefactor at \(y\),
are roots of unity by rationality of \(E_H\) and the
quadratic function defining \(\alpha\) on the Hom space.
Thus each coefficient is algebraic and of absolute
value at most one at every embedding.

Multiplication
by \(c\) sends the period lattice into the contracting
half lattice, so \(|S_B|\le c^{[E:\Q]nn'}\). The rational
lifts \(y-s\) give torsion points on the period torus.
Together with \eqref{eq:selected-kernel-regroup} this
bounds the sum of conjugate coefficient moduli in the
torsion presentation by
\begin{equation}\label{eq:selected-kernel-coeff-l1}
             \beta\,c^{2[E:\Q]nn'} .
\end{equation}
The ideal character used to define the moved vacuum
has been divided out in this presentation. The
remaining ratio of frame determinant characters has
modulus one at each embedding for the reference level.
Indeed all frame changes here are at the type primes,
where their uniformizer factors cancel against the tied
normalized local character. The surviving residue
characters have finite image. A zero entry satisfies the
bound directly.
\label{cor:selected-kernel-jets-3}
\end{proof}

\paragraph{The Gaussian bound after conjugation.}
Fix an inducing embedding \(\rho\) of the defining field.
The normalized theta vacuum in the two class lattices
is the moved vector of Section~\ref{subsec:theta-normalization}.
Here the two ranks are even, the period lattice is self-dual,
and the chosen nonzero two-torsion point survives multiplication by both
\((1+\sqrt{-d})/2\) and \((1-\sqrt{-d})/2\). The
coefficientwise conjugation of the one-variable normalized
theta and its torsion translates was described in
Section~\ref{subsec:selected-period-means}. It applies both
to the finite sum transporting the vacuum and to each
further rational translate in the indicator expansion.
Equivalently, the torsion transport acts on theta operators
with their Weil pairing.

If the phase changes by a root of unity \(\zeta\), the pointwise absolute bound is preserved, since
\[
 |\zeta a|=|\zeta|\,|a|=|a|.
\]
The induced CM type permutes the labels under embeddings
fixing the quadratic field. For an embedding exchanging
its two embeddings, also apply complex conjugation to the
coordinates. This gives entire conjugate functions in the
same labelled Euclidean metric, with the conjugate period
area \(A_\rho\), and preserves the squared average of the
vacuum. If the normalized sum is \(\sum_a c_aV_a\), its distinct
translate characters are orthogonal, also after conjugation. Thus
\[
 \left\|\sum_a\zeta_a c_aV_a\right\|^2
   =\sum_a |c_a|^2\,\|V_a\|^2\qquad(|\zeta_a|=1).
\] In particular the conjugate norm
squared equals \(N_\rho^{D_j}\) for
\(D_j=f nn'\), as computed in
Section~\ref{subsec:selected-period-means}.
\label{cor:selected-kernel-jets-4}

\begin{lemma}[Gaussian prefactor at a conjugate embedding]\label{lem:conjugate-jet-prefactor}
For the selected even ranks and a theta vacuum with self-dual
period lattice, assume its marked two-torsion point has nonzero
images under both \((1+\sqrt{-d})/2\) and \((1-\sqrt{-d})/2\).
The single-kernel Gaussian bound at an inducing embedding \(\rho\)
has a prefactor \(S\) satisfying
\[
 S^2\le N_\rho^{D_j}(8+16D_jr_E)^{D_j} e^{2(C+4)U_u}.
\]
The bound applies to either kernel at the coupled step.
\end{lemma}

\begin{proof}
Here is the pointwise bound with the covolume normalization.
Let \(r_E\) be the root discriminant. Write the period
lattice as the embedded module
\(\delta\,\mathfrak c I\operatorname{Hom}(L,L')\),
where \(\mathfrak c\) is the integral curve ideal of the
quadratic field, \(\delta\) is the period scale, and
\(I\bar I=\mathfrak D_{E/K}^{-1}\).
For a nonzero \(x\in\mathfrak c I\operatorname{Hom}(L,L')\)
the relative matrix trace \(v=\operatorname{tr}(x^*x)\)
is totally positive in the real subfield. Local frames
of the two self-dual modules show that
\(x^*x L\subset \mathfrak c\bar{\mathfrak c}I\bar I L\),
so this trace lies in
\(\mathfrak c\bar{\mathfrak c}I\bar I\).

Taking absolute ideal norms and using arithmetic-geometric
mean over the \(f\) real embeddings gives
\[
 \sum_\tau \tau(v)
    \ge f\,\mathrm N_{E/\Q}
                 (\mathfrak c\bar{\mathfrak c}I\bar I)^{1/(2f)}
    = f\,\mathrm N_{K/\Q}(\mathfrak c)\,\sqrt d/r_E .
\]
The area of the curve lattice is
\(|\delta|^2\mathrm N(\mathfrak c)\sqrt d/2\).
This gives minimum squared length at least
\(2f A_\rho/r_E\) also for the conjugate period lattice.
Its covolume is \(A_\rho^{D_j}\), by self-duality for
the imaginary part of the Hermitian form
\(H_\rho(z,z)=\|z\|^2/A_\rho\).

\proofstep{The submean estimate on the period torus.}
Let \(V\) be this conjugate theta vector or a Heisenberg
translate of it. Its weighted squared absolute value
\(|V(z)|^2 e^{-\pi H_\rho(z,z)}\) is periodic.
Translation by \(U_w\) translates this density and preserves
its mean. The ball of radius
\(R=(A_\rho/(2r_E+1))^{1/2}\) embeds in the period torus. Indeed for the shortest nonzero
period length \(\lambda_1\),
\[
 (2R)^2=\frac{4A_\rho}{2r_E+1}
    <\frac{2A_\rho}{r_E}\le\lambda_1^2 .
\]
Two points of the ball therefore cannot differ by a nonzero period.

Apply the submean inequality
at zero to \(U_z V\) on that ball, then insert the Gaussian
weight on the ball. Since its volume is
\(\pi^{D_j}R^{2D_j}/D_j!\), the squared bound at \(z\)
is
\[
 \frac{D_j!}{\pi^{D_j}}(2r_E+1)^{D_j}
    e^{\pi/(2r_E+1)}
    N_\rho^{D_j}\,e^{\pi\|z\|^2/A_\rho}
 \le 3(2D_jr_E+1)^{D_j}
    N_\rho^{D_j}\,e^{\pi\|z\|^2/A_\rho}.
\]
The last inequality uses \(D_j\ge1\) and \(r_E\ge1\).
This bound covers all classes and all the translates
in \eqref{eq:selected-kernel-coeff-l1}.

\proofstep{Summing the torsion translates.}
Apply the triangle inequality to their conjugate entire
sum. It satisfies the Gaussian hypothesis above with
\[
 S=\beta\,c^{2[E:\Q]nn'}\sqrt{3(2D_jr_E+1)^{D_j}}\,
                                  N_\rho^{D_j/2},
 \qquad r_\tau=\sqrt{\pi/A_\rho}.
\]
The Taylor differential multiplier in label \(\tau\)
is the conjugate of the generator of the coordinate axis
\(d_\tau\). The normalized metric is consequently
\(\Omega_\tau=A_\rho/(\pi|\rho(d_\tau)|^2)\),
so \(|r_\tau\rho(d_\tau)|^2=\Omega_\tau^{-1}\).
The formal jet in the coefficient field agrees with
this conjugate analytic jet.

Each relevant monomial has
the prescribed total degree at each label, and comparison
before applying \(\rho\) is by the multivariable Taylor
coefficients, multiplied by their multi-index factorials
and the labelled differential powers. The complex
coefficient form is the Gaussian pairing normalized on
the highest polynomial in each label as in
Lemma~\ref{lem:padded-jet-estimate}. Thus the
same operator bound holds on the embedded coefficient
spaces. Finally \([E:\Q]=2f\) and
\(D_j\le fm^2N_m^2\). Using \(3\le 8^{D_j}\) gives
\[
 S^2\le N_\rho^{D_j}(8+16D_jr_E)^{D_j}
                       e^{2(C+4)U_u}.
\]
This is the single-entry prefactor used in the conjugate
mean, including both entries of the last coupled step.
\label{cor:selected-kernel-jets-5}
\end{proof}

\subsubsection{The archimedean mean on the reference chain}
\label{subsec:selected-period-means}

The single-entry estimate involves both the CM period and the norm of a theta vector. The squared kernel height uses their means over the same embeddings. Write \(K=\mathbb Q(\sqrt{-d})\) for the CM curve field,
\(r_E\) for the root discriminant of the CM field \(E\), and
\(\ell_K=(L'/L)(1,\chi_{-d})\). The discriminant of \(K\) is
\(-d\), since \(d\equiv7\bmod8\) is prime. In this paragraph an
overline on a real quantity means its average over the embeddings
of a number field defining the curve, its differentials, the
finite class data and the coefficients of the selected base
sections.

The field of definition can be enlarged to contain the torsion points used
in a calculation. If \(F'/F\) is such an enlargement and \(a_\rho\) is a quantity
attached to an embedding of \(F\), then
\[
 \frac{1}{[F':\Q]}\sum_{\rho':F'\hookrightarrow\C}a_{\rho'|_F}
    =\frac{[F':F]}{[F':\Q]}\sum_{\rho:F\hookrightarrow\C}a_\rho .
\] In particular one may first average
above the infinite place of the period matrix. Every embedding
there bounds the same squared operator norm, and the two
conjugate embeddings of its totally imaginary coefficient field
give that place with the same multiplicity.

\begin{proposition}[Mean normalized differential area]\label{prop:mean-differential-area}
For the full-order CM models of \(K=\mathbb Q(\sqrt{-d})\)
with discriminant-one period lattice as in
Proposition~\ref{prop:level-free-theta-jets}, the mean normalized
area of the axis
generator \(d_\tau\) is
\[
 \overline{\log\frac{A_\rho}{\pi|\rho(d_\tau)|^2}}
   =\gamma+\log(2\pi/r_E)-\ell_K .
\]
The mean is unchanged by adjoining the finite class or
coefficient data of the selected test.
\end{proposition}

\begin{proof}
\proofstep{Average the discriminant-one area over the full ideal class group.}
Let \(A_\rho\) be the area of the period lattice at embedding
\(\rho\), with its discriminant normalized to one. The main
theorem of complex multiplication for the full order of \(K\)
makes its ideal classes equidistributed over these embeddings.
Indeed over \(K\) the Galois action on the singular moduli is
the ideal class action on the lattices, simply transitive. Complex
conjugation permutes the same classes, and adjoining the
coefficients of a discriminant-one equation gives equal
multiplicities. A homothety \(\Lambda\mapsto c\Lambda\) changes the
discriminant by \(c^{-12}\) and the area by \(|c|^2\). For two
discriminant-one representatives in the same class,
\[
 c^{-12}=1\quad\Longrightarrow\quad
 \operatorname{area}(c\Lambda)=\operatorname{area}(\Lambda).
\] \cite{cm-reciprocity}

\proofstep{Compute the average using the first limit formula.}
Choose a positively oriented basis \(\omega,\omega\tau\) in a
class, write \(y=\Im\tau\), and put
\(Z(\tau,s)=\sum' y^s/|m+n\tau|^{2s}\). Kronecker's first
limit formula gives residue \(\pi\) at one and constant term
\(2\pi(\gamma-\log 2-\log(\sqrt y\,|\eta(\tau)|^2))\).

For an ideal \(I\), its unscaled covolume is
\(\mathrm{N}(I)\sqrt d/2\). The sum over its nonzero elements
counts each integral ideal in the inverse class once for each
unit, by \(a\mapsto(a)I^{-1}\). Hence, for \(\Re s>1\),
the sum of the \(Z\)'s over the classes is
\[
       |\mathcal O_K^\times|(\sqrt d/2)^s\zeta(s)L(s,\chi_{-d}).
\]
This also identifies their Laurent coefficients at one. Divide the
constant term by the residue on both sides. The result is
\[
  2\,\overline{\log(\sqrt y\,|\eta(\tau)|^2)}
       =\gamma-\log(2\sqrt d)-\ell_K .
\]
The analytic discriminant is
\((2\pi/\omega)^{12}\eta(\tau)^{24}=1\), so
\(A_\rho=(2\pi)^2y|\eta(\tau)|^4\). We obtain
\begin{equation}\label{eq:selected-curve-area}
        \overline{\log(A_\rho/\pi)}
             =\gamma+\log(2\pi/\sqrt d)-\ell_K .
\end{equation}
\cite{kronecker-period-mean}\label{cor:selected-period-means-1}

There is a second factor from the host. With
\(\mathfrak D=\mathfrak D_{E/K}\) and the inverse half different
\(I=\mathfrak b^{-1}\), equation~\eqref{eq:host-different-factors}
gives \(\mathfrak D=\mathfrak b\bar{\mathfrak b}\) for the fixed target
character, including the primes of its conductor. Its relative
trace dual for the bilinear pairing is
\(I^*=\mathfrak D^{-1}I^{-1}=\bar{\mathfrak b}^{-1}\). In the split coordinates of
\(E\otimes_K M\), for a field \(M\) containing its embeddings,
the pairing is the sum of coordinate products. Tensor the
perfect pairing of \(I\) and \(I^*\) with \(\mathcal O_M\)
over \(\mathcal O_K\). These are projective modules, so the
pairing stays perfect.

If \(e_\tau\) is the coordinate vector on the \(\tau\)-axis,
then pairing \(ce_\tau\) with \(y\in I^*\) gives \(c\tau(y)\).
Perfectness thus identifies the ideal of allowed values as
\[
 \{c:c\tau(I^*)\subseteq\mathcal O_M\}
      =(\tau(I^*)\mathcal O_M)^{-1}
      =\tau(\bar{\mathfrak b})\mathcal O_M .
\]
It is principal for the defining field of the selected
generators \(d_\tau\). This interpretation persists after
enlargement, and descends to any defining field by faithful flatness.
Since the two conjugate factors have equal absolute ideal norm,
the ideal norm formula and the tower formula for discriminants give
\[
  \overline{\log |\rho(d_\tau)|^2}
   =\frac{\log\mathrm N_{E/\mathbb Q}\mathfrak D}{[E:\mathbb Q]}
   =\log r_E-\log\sqrt d .
\]
Subtracting from \eqref{eq:selected-curve-area} gives the
differential term with its metric and its sign,
\begin{equation}\label{eq:selected-generator-mean}
  \overline{\log\frac{A_\rho}{\pi|\rho(d_\tau)|^2}}
     =\gamma+\log(2\pi/r_E)-\ell_K=\log\Omega .
\end{equation}
This holds at each individual type label. The same calculation
at the enhanced finite primes explains why the axis generator
there has valuation zero.
\end{proof}

\paragraph{The squared norm of the theta vector.}
\begin{proposition}[Mean squared theta norm on the reference chain]\label{prop:mean-reference-theta-norm}
For the discriminant-normalized curve with its marked
nonzero two-torsion point, \(N_\rho\) is the normalized
Gaussian squared norm of its even theta section. Its mean satisfies
\[
 \overline{\log N_\rho}\le\tfrac12(\log\sqrt d-\gamma+\ell_K).
\]
The positive-part mean has the additional allowance
\(\tfrac14\log d+2\log C_0\), with \(C_0\) an absolute bound
for the even theta constants of the reduced ideal bases.
The squared norm on the standard Hom pair at step \(j\)
is \(N_\rho^{D_j}\) over each conjugate model.
\end{proposition}

\begin{proof}
\proofstep{Evaluate the Gaussian integral in the marked characteristic.}

Let \(h_\rho\) be the distinguished nonzero two-torsion point,
and let \(\theta_\rho\) be the even theta function with this
zero, normalized at the origin. In the basis above let
\(\vartheta\) denote the corresponding even Jacobi theta
constant. Its Gaussian squared norm, using normalized Haar
average, is
\[
      N_\rho=(2y)^{-1/2}|\vartheta(0,\tau)|^{-2}.
\]
Indeed integrate the squares of its Fourier series terms over
one real period. Distinct frequencies are orthogonal. For the characteristic shift \(a\in\{0,1/2\}\), integration
in the imaginary period leaves
\[
 \sum_{n\in\mathbb Z}\int_0^1e^{-2\pi y(n+a+t)^2}\,dt
    =\int_{\mathbb R}e^{-2\pi y u^2}\,du=(2y)^{-1/2}.
\] Normalization at the
origin divides by the squared theta constant.

\proofstep{Use the algebraic integer from the half-period.}

The quotient \(|\vartheta/\eta|^{12}\) is the absolute value
of a conjugate of a single nonzero algebraic integer. To verify
the assertion with this half-period, put \(e=\wp(h)\) and
\(Y=(12e^2-g_2)/4\). For the other roots \(e',e''\) of the Weierstrass cubic,
differentiating \(4(X-e)(X-e')(X-e'')=4X^3-g_2X-g_3\) at \(e\)
gives
\[
 (e-e')(e-e'')=(12e^2-g_2)/4=Y .
\] The three half-period difference formulas and
\(\theta_2\theta_3\theta_4=2\eta^3\), together with
discriminant one, give
\(
        |\rho(Y^{-3})|=|\vartheta(0,\tau)/\eta(\tau)|^{12}.
\)
For instance at \(h=\omega/2\) the two differences are
\((\pi/\omega)^2\theta_3^4,(\pi/\omega)^2\theta_4^4\)
and the characteristic is that of \(\theta_2\).

The permutations
of the roots give the other two cases with signs. In the
identity \(u=-Y^{-3}=s_i(\theta_i/\eta)^{12}\), the multipliers
are \((s_2,s_3,s_4)=(-1,1,-1)\). These are multipliers of
the complex ratios at the chosen half-periods. Each conjugate satisfies the monic integral equation over
\(\mathbb Z[j]\)
\[
 (u-16)^3-ju=u^3-48u^2+(768-j)u-4096=0 .
\] This follows on substituting
\(\lambda=\theta_2^4/\theta_3^4\) in
\(j=256(1-\lambda+\lambda^2)^3/(\lambda^2(1-\lambda)^2)\).
Thus \(Y^{-3}\) is integral since \(j\) is an algebraic integer.
Both \(e\) and the discriminant-one model conjugate with the
specified two-torsion datum. The product formula gives
\(\overline{\log|\vartheta|}\ge\overline{\log|\eta|}\).

Combining the Gaussian identity with the area formula yields
\[
        \overline{\log N_\rho}
        \le\tfrac12(\log\sqrt d-\gamma+\ell_K).
\]
A reduced ideal basis gives \(1/2\le y\le\sqrt d/2\)
by the reduced binary quadratic form of discriminant \(-d\).
The defining Gaussian series bounds each even theta constant
at zero by an absolute \(C_0\ge1\) on this range. Thus the negative part obeys the pointwise estimate
\[
 (-\log N_\rho)_+
 =\big(\tfrac12\log(2y)+2\log|\vartheta(0,\tau)|\big)_+
 \le \tfrac14\log d+2\log C_0 .
\]
Adding its mean to \(\overline{\log N_\rho}\) bounds the mean
of \((\log N_\rho)_+\). This permits a coefficient whose
sign can vary elsewhere in the height comparison.
\cite{elliptic-constants}\label{cor:selected-period-means-2}

\proofstep{Transport the normalized one-variable section.}

For the standard source and target the rank of the quadratic
CM Hom space is \(D_j=f N_{j-1}N_j\). We spell out the use of
the theta normalization under conjugation. On its fourfold sum
the normalized vector of Section~\ref{subsec:theta-normalization}
is a normalized finite sum of Heisenberg translates of the
product of \(4D_j\) one-variable sections. At the conjugate
curve the one-variable factor is again the even section
normalized at the origin with zero \(h_\rho\). Let \(s\) denote this normalized section in its local coordinate. Evenness and the logarithmic derivative formula give
\[
 s'(0)=0,\qquad (\log s)''(z)=-\wp(z+h_\rho)-s_2 .
\]
Formula~\eqref{eq:cm-isogeny} computes \(s_2\) algebraically from the nonzero points of a CM isogeny with \(a^2\ne a\bar a\). These identities conjugate
coefficientwise and determine the entire Taylor series.

\proofstep{Compare translated sections and the fourth tensor.}

Write \(s\) for the unshifted section and \(s_t\) for its
translate by a point \(t\) of order \(m\). The quotient has rational
\(m\)-th power, and
\[
            \operatorname{div}\bigl((s_t/s)^m\bigr)
                      =m([h-t]-[h]).
\] After conjugation it is a scalar times the
corresponding rational function on the conjugate curve, by
equality of divisors. More precisely the functions can be taken
over the algebraic coefficients. Their local Laurent series is
algebraic coefficientwise, and the space of rational functions
with that divisor is a line defined over the field of the torsion
points. The product of this rational quotient over the \(m\)
translations by \(jt\) is a root of unity by the period and
commutator laws, as in Section~\ref{subsec:baseline-jets}.

If \(f_1=c f_2\) are the two rational quotients with the
same divisor after conjugation, their translated products differ by
\[
 \prod_{j=0}^{m-1} f_1(z+jt)
     =c^m\prod_{j=0}^{m-1} f_2(z+jt).
\]
Both products are roots of unity in this normalization, hence
\(|c|^m=1\) at every embedding.
Taking the formal \(m\)-th root gives the translated conjugate
section up to a unit of complex modulus, including a translate
vanishing at zero. On the finite quotient, let \(\kappa(x,a)\) be the commutator
character of the translated summand \(G_a\) under the unitary
translation \(W_x\), with \(a\) the corresponding torsion class.
For \(a\ne b\), perfection of the pairing gives \(x\) such that
\(\kappa(x,a)\ne\kappa(x,b)\). Thus
\[
 \langle G_a,G_b\rangle
  = \kappa(x,a)\overline{\kappa(x,b)}
       \langle G_a,G_b\rangle=0 .
\]
The transported classes remain distinct and their commutator
pairing is the conjugate Weil pairing. The same calculation
therefore gives orthogonality after conjugation.

Changing the
phases of the individual summands preserves the norm. The rational CM isometry conjugates to an
isometry, since its adjoint equations are algebraic over the
CM fields, on which conjugation commutes with the embeddings.
Its change of variables preserves normalized averages.
For the index \(I\) of this theta-line transfer the conjugate
section is, up to the determinant scalar of modulus one in
\eqref{eq:kappa-determinant}, \(I^{-1/2}\sum_{a=1}^I\zeta_a G_a\).
Each \(|\zeta_a|=1\), and the \(G_a\) are orthogonal with
squared norm \(N_\rho^{4D_j}\) since each is a unitary translate
of that one product section. Hence
\[
      \left\| I^{-1/2}\sum_{a=1}^I\zeta_a G_a\right\|^2
               = I^{-1}\sum_{a=1}^I\|G_a\|^2
               = N_\rho^{4D_j}.
\] Taking fourth roots of squared
tensor norms now gives \(N_\rho^{D_j}\) for the squared norm
of the standard vector. The ideal-character multiplier on the
standard pair is one.

\end{proof}

\paragraph{Averaging the selected error from the kernel.}
\begin{proposition}[A uniform bound for the mean logarithmic lift norm]\label{prop:mean-reference-lift-price}
For the selected word and reference chain, write \(q_{jz}=\sum_i\nu_{jzi}\) for the size of the partition at lift \(j\) and real place \(z\). Set \(\varphi(u)=u(\log u-1)\) for \(u>0\) and \(\varphi(0)=0\). The mean logarithmic lift bound \(h_B(j)\) satisfies
\[
 h_B(j)\le
 \sum_z\left(\sum_i\varphi(\nu_{jzi})-q_{jz}\log\Omega\right)
       +C fN_m^2(L_a+m^3L).
\]
The constant is independent of the CM field, word, embedding and lift index \(j\)
for fixed \(C_1\), and the estimate applies at both coupled
indices with the given normalization.
\end{proposition}

\begin{proof}
\proofstep{Start from the squared entry bound.}

Put \(n=N_{j-1}, n'=N_j, v=r_E/2\), \(L_a=\log(a r_E)\),
\(q_{jz}=\sum_{i=1}^n\nu_{jzi}\) and
\(\varphi(s)=s(\log s-1)\), with \(\varphi(0)=0\).
For each embedding inducing a given infinite place of the
coefficient field, the single entry in that place has the
logarithmic squared bound
\begin{align*}
 k_j(\rho)={}&\sum_z\left(\sum_i\varphi(\nu_{jzi})
       -q_{jz}\log\frac{A_\rho}{\pi|\rho(d_z)|^2}\right)
      + D_j\log N_\rho + D_j\log(8+32D_jv)\\
 &\quad +2C_1 U_u+
           \sum_z\bigl(1+2nn'\log(2q_{jz}+2)\bigr).
\end{align*}
Here \(C_1=C+4\) by the pointwise estimate in
Section~\ref{subsec:selected-kernel-jets}.
The ordinary and coupled sums now use the multiplicities just
proved. Equivalently one may put \(C_1+58\) in the last bound
for each lift. The upper bound for each lift also includes the
nonnegative logarithms of its two class and two index counts.
These add at most \(112U_u\) per lift if retained explicitly.

\proofstep{Insert the parameter bounds and the two archimedean means.}

Since \(\log g_i\le9L+1\), \(\log T_u\le25mL\) and
\(L,m\ge1\), we have \(U_u\le45fm^3N_m^2 L\).
Indeed,
\[
 \log(g_0g_1T_u)\le18L+2+25mL\le45mL .
\]
Moreover \(fN_m\le b\), \(q_{jz}\le3aN_m\),
\(4x\le b\), \(a=xb\) and \(r_E\ge1\). These give
\(\log(8+32D_jv)=O(L_a)\) and
\(\log(2q_{jz}+2)=O(L_a)\) with absolute constants.

The selected bounds on \(\log d\) and \(|\ell_K|\) are at
most \(L\). Applying the two archimedean means above shows
that the averaged lift price is at most
\[
  \sum_z\left(\sum_i\varphi(\nu_{jzi})-q_{jz}\log\Omega\right)
       +C fN_m^2(L_a+m^3L).
\]
This is the numerical height contribution \(h_B(j)\) on the
same word and at both coupled indices, with \(C\) independent
of the host, word, embedding, and \(j\), for fixed \(C_1\).
The terminal evaluation follows in the next subsection. The
row scalar enters the comparison with the word height.
\label{cor:selected-period-means-3}
\end{proof}

\subsubsection{The height bound for the terminal evaluation}
\label{subsec:selected-terminal-height}

At the compact end of the chain the height also includes the branching evaluation. Its two interlacing weights determine which coefficients have to be bounded.

Use the reference sequence and the terminal point and row spaces at the refined levels of Section~\ref{subsec:terminal-realization}. At each
real label \(z\), the terminal large and small weights
\(\lambda,\sigma\) have ranks \(N=N_m\) and \(N-1\).
They interlace and have the same last entry. Put
\(A_i=\sigma_i-\lambda_{i+1}\) and \(B_i=\lambda_i-\sigma_i\)
for \(1\le i<N\), and let
\[
 D_{0,z}=\sum_i A_i,\qquad
 n_z=\sum_i\lambda_i-\sum_i\sigma_i=\lambda_1-D_{0,z}.
\]
The harmonic weight on the last diagonal coordinate is
\(n_z=\lambda_N+\sum_i B_i\). In this equality the last
entry of the full weight accounts for the determinant
character in the polynomial model.

\begin{proposition}[Height of the terminal evaluation]\label{prop:terminal-evaluation-height}
For the reference sequence, let \(T\) inside the sum below
denote the selected top labels and \(T_u\) the assigned prime.
The logarithmic squared norm of terminal evaluation in the
indicated point metrics is bounded above by
\[
 -\sum_{z\in T}a_z\log T_u
 +O\!\left(L_a\bigl[|T|b'N+
     f(S_c+m\eta)N+fN^2\bigr]\right).
\]
The implicit constant is uniform for the selected ranges of the parameters.
\end{proposition}

\begin{proof}
\proofstep{Separate the compact branching norm from its Capelli scalar.}
Recall the branching embedding \(J=c_zJ_0\), where
\(J_0h_\sigma=\xi_z\) and \(c_z\) is the nonzero integer
obtained from the capped interlacing intervals. With the
invariant metrics in which both highest vectors have norm
one, \(J_0^*J_0=\|\xi_z\|^2\,1\) by Schur's lemma.
Since \(J=c_zJ_0\), Schur's scalar for the normalized map gives the squared norm of the normalized transpose:
\[
 |c_z|^2\|\xi_z\|^2 .
\] The branching polynomial \(\xi_z\)
is a product of the initial-column flag minors in the
smaller-rank model inside the large one. Its normalized
squared norm is at most one. Indeed realize these minors
as coefficients in tensor powers of exterior powers, with
the wedge basis orthonormal. The tensor of the initial
wedges spans the highest weight of the Cartan summand.

Let \(u\) be the tensor of row wedges, of norm one, and
\(P_{\rm Car}\) the orthogonal projection to the Cartan summand.
The coefficient vector representing \(\xi_z\) is \(P_{\rm Car}u\), so
\[
 \|\xi_z\|^2=\|P_{\rm Car}u\|^2
             =1-\|(I-P_{\rm Car})u\|^2\le1 .
\]
The unit of comparison is the initial wedge tensor
itself, which gives \(h_\lambda\). In particular the
weaker bound \(\log\|\xi_z\|^2\le N(N-1)\log(2a)/2\)
follows since \(a\ge1\) and \(N\ge1\).

For the Capelli scalar, let
\(D_{1,z}=\sum_i\min(\lambda_i-\lambda_{i+1},D_{0,z})\).
This is exactly the number of trial integers in the capped
intervals after omitting the selected \(\sigma_i\).
Each trial contributes \(N-1\) differences. Since all shifted
weights and trial integers lie in the common spread interval,
each difference has absolute value at most \(2a\), by the
selected spread bound. Consequently
\[
             \log |c_z|\le (N-1)D_{1,z}\log(2a).
\]
\label{cor:selected-terminal-height-1}

\proofstep{Use the signed degree in the common class metric.}
At each class the last line is scaled by the fixed word
element \(\pi_u\), with \(\pi_u\overline{\pi_u}=T_u\).
The action in the transported evaluation includes
\(\prod_z z(\pi_u)^{-n_z}\) with the full signed \(n_z\).
Its squared modulus is \(T_u^{-\sum_z n_z}\).
Both metrics in the terminal pairing use the same row
class index and the same reciprocal stabilizer weights.
In particular the large fibre at a row class is the
terminal point fibre prescribed over that class. Applying
Cauchy--Schwarz with these weights gives the squared
pairing bound by
\[
 T_u^{-\sum_z n_z}\prod_z |c_z|^2\|\xi_z\|^2
\]
times the squared norms in the two terminal point metrics.
This estimate starts after the coupled map of
Section~\ref{subsec:selected-norms}; its two-kernel
class sum is already in that map. The evaluation and
quotient-metric comparisons in
Section~\ref{subsec:selected-lift-scalar} give its
bound on the framed spectral row.
\label{cor:selected-terminal-height-2}

\proofstep{Bound the head and the tail by interlacing.}

Write the head Harish--Chandra coordinates at a label as
\(u_0,\ldots,u_{k-1}\), with the interlacing coordinates with smaller weights
\(v_0,\ldots,v_{k-2}\). Write \(b_\alpha,e_\alpha\),
\(0\le\alpha<m\), for the two ends of the successive
tail blocks of size \(t+1\). The coordinates with smaller weights lie
in \([0,S_c]\), \(u_1<v_0\), and
\(|b_\alpha-e_\gamma|\le\eta\). With zero-based indices the head entries in the large and
small weights are, respectively,
\[
 u_i-\frac{k-1}{2}+i-\frac{m(t+1)}2,\qquad
 v_i-\frac{k-2}{2}+i-\frac{m(t+1)}2 .
\]
Their difference on the shared indices is \(u_i-v_i-1/2\).
The subsequent block \(\alpha\) in the large weight is constant at
\(-e_\alpha+k/2+(2\alpha+1-m)(t+1)/2\).

In each tail block the \(t\) later entries agree with
the corresponding entries of the large weight in the
shifted index for \(A_i\). The first entry in this tail block contributes the difference
\[
 e_\alpha-b_\alpha-\tfrac{t+1}{2}
\]
to \(D_{0,z}\).
This follows by substituting the first coordinate
\(-b_\alpha\) and the \(t\) successive coordinates
from \(-e_\alpha\) in the shifted Harish--Chandra row
of the small representation.
\label{cor:selected-terminal-height-3}

The head sum of \(A_i\) is at most
\(v_0-u_{k-1}-k+3/2\) by telescoping with interlacing,
and the tail contributes at most \(m\eta\).
Off the top labels \(u_{k-1}=v_{k-2}-1/2\).
At a top label \(z\), use
\(b_z=-u_{k-1}-(k-1)/2\le2b'\).
It follows that
\[
        D_{0,z}\le S_c+m\eta+2b'\,{\bf1}_T(z),
\]
where \(T\) is the set of top labels.

To bound \(D_{1,z}\),
cap the first gap and the gap between the head and the
first tail block by \(D_{0,z}\). The head gaps left
over telescope from the second to the \(k\)-th entry.
Their sum is bounded by \(D_{0,z}+S_c\): subtract the
single term comparing \(v_{k-2}\) with the last head
entry and use \(u_1<v_0\), \(v_{k-2}\ge0\).
The remaining tail gaps telescope to
\(e_{m-1}-e_0-(m-1)(t+1)\le\eta\), since
\(b_0\le e_0\) and \(|b_0-e_{m-1}|\le\eta\).
Combining the two capped gaps, the remaining head gaps and
the telescoping tail gives, respectively,
\[
 D_{1,z}\le2D_{0,z}+(D_{0,z}+S_c)+\eta
                  =3D_{0,z}+S_c+\eta .
\]

At a top label \(\lambda_1=a_z-(N-k)/2\).
Otherwise \(u_0-v_0\ge S_c\) and \(v_0\ge0\)
give \(\lambda_1\ge S_c-(N-1)/2\), which is sufficient
even when \(n_z\) is negative. Therefore
\[
 \sum_z n_z\ge \sum_{z\in T}a_z-2|T|b'-f(S_c+m\eta+N).
\]
Sum the logarithmic bounds for the branching scalars.
Using \(\log(2a)=O(L_a)\) and
\(\log T_u\le25mL\le25N L_a\) gives the logarithmic
squared evaluation price
\[
 -\sum_{z\in T}a_z\log T_u
 +O\!\left(L_a\bigl[|T|b'N
        + f(S_c+m\eta)N+fN^2\bigr]\right).
\]
This is \(h_A\) with an absolute uniform error constant,
on the same reference sequence as the \(h_B(j)\) bounds.
\label{cor:selected-terminal-height-4}
\end{proof}


\subsubsection{The lift-adjoint scalar in the selected row metric}
\label{subsec:selected-lift-scalar}

With the kernel norms and terminal evaluation in hand, the remaining normalization is the scalar of a lift followed by its adjoint in the actual row metric.

Fix a lift \(j\) at the reference tuple. In the adelic
calculation write \(F^+\) for the real field, \(E\) for
its CM extension, and \(f=[F^+:\Q]\). The row rank \(r\), definite rank \(n_V\) and stable-range gap satisfy
\[
             r=k+m-1,\qquad n_V=r+t,\qquad t\ge r+5 .
\] The definite source has the
standard self-dual lattice with an integral orthonormal
frame. The local level and Weil vectors use the specified
split half places, with at most one place above each prime
of \(F^+\). The characters intertwine the same local
level actions on the vectors and the forms, and the
Schwartz vector is spherical away from those places.
We use the level, weight, rational Gram and character
identifications of the selected primitive rows. These
conditions apply at the reference tuple, including its
conjugate coefficient realizations.
\label{cor:selected-lift-scalar-1}

\begin{proposition}[Scalar bounds and reference height]\label{prop:selected-reference-height}
For a selected reference sequence \(u\), let \(B_u\) be
the comparison operator in the framed spectral row. With
the common row metric and the period shift
\(\bar\kappa=\overline{\log\kappa}\), its height satisfies
\[
 \operatorname{ht}(B_u)\le h_A(u)+
       \sum_{j=1}^m\bigl(h_B(j)-h_X(j)\bigr)-\bar\kappa .
\]
Here \(h_A\) is the terminal bound, \(h_B(j)\) bounds the
kernel height at lift \(j\), and \(h_X(j)\) bounds the
logarithm of the adjoint scalar \eqref{eq:selected-adelic-scalar}
on the occurring primitive rows. The estimates are uniform on the reference
sequence data, including the coupled end steps.
\end{proposition}

\begin{proof}
\proofstep{Unfold the adjoint in the weighted form and multiplicity metrics.}
Let \(T\) be the full theta lift from the definite form
space with its weighted class metric to the cusp
multiplicity space \(H_W\). This space has the row
metric, which is \(\prod_z d_z\) times the Petersson
multiplicity product, where \(d_z\) is the dimension
of the lowest compact type at \(z\). This metric can
be computed with a unit vector of that type at each
real place by Schur orthogonality. For the finite level group \(P\), the normalized and
unnormalized adjoints on the stated metrics are related by
\[
 T^\dagger=|P|^{-1}T^*,\qquad
 \langle T^\dagger U,T^\dagger U'\rangle
    = |P|^{-2}\langle T^*U,T^*U'\rangle .
\]
For the reference scalar \(|P|=1\).

Let \(U\) lie in a nonzero simultaneous off-level Hecke
part and let \(U'\in H_W\). Let \(C\) be the \(d_C\)-dimensional coefficient module
on the definite side, with orthonormal basis \((e_\nu)\).
Write \(\mathcal R_\nu(U)\) for the scalar reverse
theta integral against \(e_\nu\) in the Petersson
normalization. The source-level transformation law makes
\((\mathcal R_\nu(U))_\nu\) a definite form of the required
character. Put \(D_\infty=\prod_z d_z\). The row metric
contributes \(D_\infty\) to each adjoint, so the weighted
class pairing is
\[
 \langle T^*U,T^*U'\rangle
   =D_\infty^2\sum_{\nu=1}^{d_C}
      \langle\mathcal R_\nu(U),\mathcal R_\nu(U')\rangle_{\rm Pet,def}.
\]
Here the definite Petersson pairing is the sum over its
classes weighted by reciprocal stabilizer orders. The definite quotient is compact,
and cuspidal rapid decay gives the integrability on
the other side.

The convergent Rallis inner-product identity for these
reverse integrals unfolds the product into the finite
oscillator matrix coefficient and the real Fock and
discrete-series coefficients. This is the unitary
doubling identity for cuspidal forms, with the
anisotropic definite member and with the matching
flat section given by the two Schwartz vectors. The
shift is \(s_0=t/2\). In Yamana's Theorem 2, the
doubled skew-Hermitian space has dimension \(2r\),
the definite Hermitian space has dimension \(n_V\),
and its rational Witt index is zero. The range in Yamana's theorem puts the special point in absolute convergence for the Eisenstein series, since
\[
 n_V>2r\quad\Longrightarrow\quad (n_V-r)/2>r/2 .
\] The
Eisenstein series equals the theta integral with
coefficient one for the quotient measure of mass one.

Write \(d\mu_1\) for the definite probability measure and
\(d\mu\) for the definite class-sum measure on the quotient \(Y\).
The conversion in the theta integral is
\[
   \int_Y \theta\,d\mu=\mu(Y)\int_Y\theta\,d\mu_1 .
\] Unfolding and factoring the spherical zeta
integrals gives the partial standard factor at
\(s_0+1/2\). The retained local integral is absolutely
convergent there. The finite zeta evaluation and the two
real evaluations below use the same covolume-one
oscillator measure and real group measures.
\cite{yamana-siegel-weil}\label{cor:selected-lift-scalar-2}

Let \(v_z\) be a unit lowest-type vector and \(H_z^+,H_z^-\) the hook products of the positive and negative partitions of the actual Fock vector at \(z\). Its matrix-coefficient factor is
\[
 J_z=\int |\langle v_z,\pi_z(g)v_z\rangle|^2\,dg .
\] Let \(I_V\) be the index of the source
level in its self-dual stabilizer, \(M_V\) the mass at
that stabilizer, and \(N_\Phi\) the squared norm of
the finite Weil vector.

Write \(S\) for the finite
places omitted in the spherical product, \(D_S\)
and \(D^S\) for the two parts of the normalizing
Euler product, and \(L^S\) for the standard factor
including its unitary character twist. The unfolding
of one reverse integral yields the factor
\[
 I_VM_V\Bigl(\prod_z H_z^+H_z^-J_z\Bigr)
       |P|\,N_\Phi\, L^S/D^S
\]
times the Petersson multiplicity product. The
level factor \(|P|\) counts the classes in the
finite sum that differ by the level group, on which
the two characters cancel. The type determines
\(L^S\) through the spherical eigenvalues of the
first argument. Taking its diagonal entry for a
nonzero \(U\) shows that the resulting scalar is
real and nonnegative. The positive lower estimates
below show it is positive.

In particular the
absolute value of the standard factor equals its
value in that diagonal formula. Let \(A_{\rm rev}\)
denote the displayed factor for one normalized coefficient
basis vector. The reverse products carry \(D_\infty^2\)
from the two adjoints. Since
\(\langle U,U'\rangle_{\rm row}
      =D_\infty\langle U,U'\rangle_{\rm Pet}\),
the sum over the orthonormal coefficient basis is
\[
 \langle T^*U,T^*U'\rangle
     =D_\infty\,d_C A_{\rm rev}\langle U,U'\rangle_{\rm row}.
\]
Here \(A_{\rm rev}\) includes the counted factor \(|P|\). Dividing by \(|P|\)
gives, for every \(U'\),
\(\langle TT^\dagger U,U'\rangle=\lambda\langle U,U'\rangle\),
where
\begin{equation}\label{eq:selected-adelic-scalar}
 \lambda=\Bigl(\prod_z d_z\Bigr)d_C M_V
       \Bigl(\prod_z H_z^+H_z^-J_z\Bigr)
       \frac{|L^S|}{D_SD^S}.
\end{equation}
Here the exact finite-level identity used in the last
line is \(I_V N_\Phi=D_S^{-1}\).
\label{cor:selected-lift-scalar-3}

\proofstep{Bound the finite Euler, level, and stabilizer factors.}

At a split place of residue size \(q\), reducing the
source level gives a parabolic fixing the specified
first \(r\) columns up to a subgroup \(P_q\) of
\(\mathrm{GL}_r(\mathbf F_q)\). Its last \(t\times t\)
block is arbitrary invertible, with \(rt\) free
off-diagonal entries. The relative index and
the squared oscillator norm there are
\[
 \frac{|\mathrm{GL}_{r+t}(\mathbf F_q)|}
      {|\mathrm{GL}_t(\mathbf F_q)|q^{rt}|P_q|},
 \qquad |P_q|q^{-r(r+t)}.
\]
The second equality counts admissible \(r\)-column
matrices, each with a unit-modulus coefficient, in
the integral matrix ball. Of \(r\) columns in dimension \(t+r\), the choices for column
\(j+1\) after \(j\) independent choices are \(q^{t+r}-q^j\).
Their proportion in the full residue matrix space is
\[
 \prod_{j=0}^{r-1}\frac{q^{t+r}-q^j}{q^{t+r}}
       =\prod_{a=1}^r(1-q^{-t-a}),
\]
the reciprocal local normalizer.

At the other places the vector is the
vacuum of squared norm one and the source has its
self-dual level. The split half-place condition allows
multiplication over distinct \(F^+\)-places. This
proves the finite-level identity, also with prescribed
local character factors. The stabilizer index and mass use
the same measure, so the quotient volume at the level
is \(I_VM_V\).

For \(w\) over a spherical place let \(q_w=\mathrm Nw\).
The unramified zeta integral gives an inverse Euler
polynomial whose \(r\) parameters have modulus at most
\(q_w^{(r-1)/2}\) in the normalization before evaluating
at \(s_0+1/2\). This bound uses the spherical part of the
unitarizable cuspidal representation. Consequently
each factor of \(L^S\) is \((1-x_{wa})^{-1}\) with
\[
 |x_{wa}|\le q_w^{(r-1)/2-(1+t)/2}\le q_w^{-3}.
\]
There are at most \(2f\) primes over a rational
prime and \(\sum_{n\ge2}n^{-3}<1\). Thus the Euler
product converges absolutely, and
\(-\log|1-x|\ge-|x|\) gives
\[
 \log|L^S|\ge-\sum_{w\notin S}\sum_{a=1}^r q_w^{-3}
           \ge-2fr\sum_{n\ge2}n^{-3}\ge-2fr .
\]

At a base place of norm \(q\)
the normalizer is
\(\prod_{a=t+1}^{t+r}(1-\epsilon_q^a q^{-a})^{-1}\)
with \(\epsilon_q\) the local quadratic sign in the
doubling normalization. The \(r\) arguments have
absolute value at most \(q^{-2}\). The inequality \(-\log(1-q^{-2})\le2q^{-2}\) and the bound of \(f\) places over each rational prime give, for either positive product \(D\in\{D_S,D^S\}\),
\[
 \log D\le 2fr\sum_p p^{-2}\le4fr .
\] These inequalities are independent of
the finite set omitted and of the occurring Hecke type.

Let \(w_E\) be the number of roots of unity of \(E\).
The identity class alone gives
\(M_V\ge(w_E^{n_V}n_V!)^{-1}\). Indeed an integral
unitary matrix in the standard frame has one
nonzero entry in each column. For a nonzero integral
entry, the mean squared modulus over the complex
embeddings is at least one by the norm and the
arithmetic--geometric mean inequality. Indeed for a column \((a_i)_i\) the relation and positivity give
\[
 1=\sum_i\frac1{[E:\Q]}\sum_\sigma|\sigma(a_i)|^2
   \ge\sum_{i:a_i\ne0}|N_{E/\Q}(a_i)|^{2/[E:\Q]}
   \ge\#\{i:a_i\ne0\}.
\]

The unit length forces exactly one nonzero entry, and equality in
the mean inequality forces all of its conjugates to have the same
absolute value. Thus the one entry has
modulus one at every embedding, so it is a root
of unity. The same argument applies to all columns,
and invertibility forces the positions to form a
permutation. Euler's totient bound \(\varphi(a)^2\ge a/2\), applied to
the roots of unity in \(E\), now gives
\[
              w_E\le 2\varphi(w_E)^2\le 8f^2,\qquad
              M_V\ge\frac{1}{(8f^2)^{n_V}n_V!}.
\]
\label{cor:selected-lift-scalar-4}

\proofstep{Retain the compact hook factors in the lower bound.}

The lowest polynomial is a product of initial minors
in the positive block and reverse initial minors in
the negative block, with exponents successive
differences of the two partitions. The two blocks
use disjoint sets of variables. In covolume-one
Gaussian coordinates the squared Fock norm is
therefore the product of the two hook products.
Each block uses \eqref{eq:selected-fock-matrix}
on its highest matrix coefficient. A determinant
character changes neither its compact normalization
nor the invariant coefficient pairing. This gives
the \(H_z^+H_z^-\) factor in
\eqref{eq:selected-adelic-scalar}.

For the regular half-integral Harish--Chandra
parameter \(\gamma_{z,1}>\cdots>\gamma_{z,r}\)
of the selected discrete series, Schur orthogonality
and the formal-degree formula give
\[
 J_z^{-1}=\frac{1}{V_r}
       \prod_{a<b}\frac{|\gamma_{z,a}-\gamma_{z,b}|}{b-a},
 \qquad
 V_r=\prod_{\ell=1}^r\frac{2\pi^\ell}{(\ell-1)!}.
\]
For compact probability measure the density at the
identity, in the comparison Lie algebra coordinates,
is \(2^{-r(r-1)/2}/V_r\). Indeed the quotient of
\(U(\ell)\) by \(U(\ell-1)\) contributes the unit
sphere of real dimension \(2\ell-1\), with volume
\(2\pi^\ell/(\ell-1)!\) in the column coordinates.
The paired off-diagonal coordinates in the chosen
Lie algebra volume supply the factor \(2^{-r(r-1)/2}\)
in the probability density.

The compact-dual measure
matches that density on \(U(p_z,q_z)\). Its inverse
coefficient integral is the displayed root product
without \(V_r^{-1}\). The density ratio at the identity is
\[
 \frac{2^{-r(r-1)/2}}{2^{-r(r-1)/2}/V_r}=V_r .
\]
Multiplication by this constant converts the compact-dual
coefficient integral to the real measure in the unfolding. The
strict parameter inequalities give integrability and
positivity of \(J_z\). If \(d\) is the degree of the
local Fock polynomial, the parameter width is at
most \(d+n_V+3\), by the lowest-type shift. Each
denominator \(b-a\) is at least one. The volume
formula gives \(V_r\ge r^{-r^2}\), so
\(\log J_z^{-1}\le r^2\log(r(d+n_V+3))\).
\cite{harish-chandra-degree}\label{cor:selected-lift-scalar-5}

Put \(L'=\log(b\,\mathrm{rd}(E)\,g_0g_1 n_V)\).
At the selected sizes,
\(L'\ge5.5\), \(\log n_V\le L'\), and the degree
bound \(d\le3rb\) gives \(\log(d+2)\le2L'\).
Because \(n_V\ge2r\) and \(n_V\ge8\),
\[
 \log\big(r(d+n_V+3)\big)\le\log r+\log n_V+\log(d+2)
    \le4L' .
\]
Summing the real loss over the \(f\) places costs at most
\(4fr^2L'\). The mass loss is at most
\[
 n_V\bigl(4f+\log 2+L'\bigr)
\]
by \(\log w_E\le4f+\log 2\).
The finite loss from the Euler, level and stabilizer factors
of \eqref{eq:selected-adelic-scalar} is at most \(10fr\).
Both dimension factors in \eqref{eq:selected-adelic-scalar}
are positive integers. Adding the losses gives
\begin{equation}\label{eq:selected-hook-lower}
       \log\lambda\ge
        \sum_z\log(H_z^+ H_z^-)-8 f n_V^2 L'.
\end{equation}
In particular the error bound is the same for every
occurring type, unit vector and embedding at a lift.

\proofstep{Descend the scalar at each conjugate embedding of an occurring part.}

At the base embedding the untwisted principal-level
parts span the type decomposition of the full multiplicity
space. They are orthogonal for the row product. The good
type projector \(Q\) is orthogonal and commutes with
this decomposition. The actual lift is \(S_0=QT\).
On a type on which \(TT^*=\lambda I\), the compressed operator
satisfies
\[
 S_0S_0^*=QTT^*Q=\lambda Q .
\]
Thus \eqref{eq:selected-hook-lower} holds on every nonzero
type in the projected image.

At another inducing embedding, use the rational lift
with the selected phase one, its conjugate splitting
characters and its conjugate finite vector. The
conjugate Gram identification gives an isometry of the
definite source forms for the form obtained by applying
that embedding to the rational source pairing. For an
embedding fixing the inducing quadratic field, express the rational form on the definite classes as a sum of geometric theta evaluations with algebraic coefficients.
The same class sum with embedded coefficients evaluates
the new integral, by the conjugate coefficient map and
the covariance of theta evaluation.

Let \({\rm ev}\) denote evaluation on the chosen extremal vector. Equality after this evaluation suffices in the multiplicity space because
\[
 {\rm ev}(v)={\rm ev}(v')\quad\Longrightarrow\quad v-v'\in\ker({\rm ev})=0 .
\] This identifies
the good projection of the new lift with its rational realization up to the unitary change of row coordinates, the action of
a central level element, and a unit-modulus scalar.
The last factor is a product of central character
values, a root of unity, and the unit phase of the
theta comparison. For this unitary map and the central action write their
product as \(U\) and the remaining scalar as \(c\). Relative to
the normalized row Gram matrix \(G\) this gives
\[
 (cU)^*G(cU)=|c|^2 U^*GU=G .
\] At the remaining embeddings
use complex conjugation of both Gram pairings. The
Petersson product is computed anew in the conjugate
model.

For a retained type, take the source vector \(x\)
in the tied rational adjoint image. At the base
embedding its Gram pairing with the rational kernel
is zero. This vanishing is a rational class-sum identity,
so holds at the other embeddings as well. The conjugate
source realization is an isometry, and the rational kernel
spans the Hilbert-model kernel after scalar extension.
Thus the realized source vector is in the orthogonal
complement of the conjugate lift's kernel. The primitive
type and lift identifications give, in its metrics,
\[
 \|S_v x\|_v^2=\lambda_{j,p}(v)\,M_v(x,x).
\]
Here \(M_v\) is the source Gram majorant and the
conjugate realization of the target uses the same
row basis as in the period identity. Both products
carry the reciprocal level index.

In the definite
Hilbert model the majorant is the full coefficient
product divided by that index. The conjugate
target norm is divided by the same index. To obtain
the scalar relation one uses the conjugate
class-sum eigenspace identity on the rational lift and
the conjugate Gram identity on the source. This
identifies the spectral labels before comparing
norms. On \((\ker S_v)^\perp\) the squared singular
values of \(S_v\) are the positive eigenvalues
of \(S_v S_v^*=QTT^*Q\). If \(s_{\min}\) is the smallest positive singular value on the
occurring part, expand a vector \(x\) in its orthonormal right
singular basis. This gives
\[
 \|S_vx\|_v^2=\sum_a s_a^2|x_a|^2
       \ge s_{\min}^2\sum_a|x_a|^2=s_{\min}^2 M_v(x,x).
\] A type present in the
coefficient comparison has some \(x\) with
\(M_v(x,x)>0\), which allows cancellation in the
scalar relation.

Conjugation permutes the labelled weight data in the
hook sum. The negative hooks are positive integers
and may be dropped. The remaining hooks are
\(H_{jz}\) of the fixed positive row partitions.
Consequently for \(C_{\rm sz}\ge8\) a uniform logarithmic lower bound is
\begin{equation}\label{eq:selected-hX}
 h_X(j)=\sum_z\log H_{jz}-C_{\rm sz} f n_V^2 L',
 \qquad \log\lambda_{j,p}(v)\ge h_X(j)
\end{equation}
on all occurring types and at all inducing embeddings.
A type with zero part makes no contribution to the
norm comparison. If a scalar is assigned to that
label in the numerical record, it is assigned the
exponential of the stipulated logarithmic lower bound, as in the
absent-part clause of the scalar record.
\label{cor:selected-lift-scalar-6}

\proofstep{Divide by the lift scalars and add the terminal cost.}

The terminal pairing first uses the weighted class metrics
of Proposition~\ref{prop:terminal-evaluation-height}. At the
same reference level \(P_0\), let \(V\) be the inclusion of
the good row basis and let \(H\) be the Gram matrix of the
principal class sum. With all matrices evaluated at the place in question, the framed norm of a test \(F\) is given by
\[
 \|F\|^2=\frac{\bar F^{t}V^*HV F}{|P_0|}.
\] The principal-level
restriction of the theta evaluation has coordinates \(VF\).
In fact this restriction on full-level forms is the level
idempotent. It cancels the section and cutoff in the
evaluation, since \(V\) already lies in the good level
image. The matrices on both sides are evaluated in the
same basis. The finite-level companion compares the two evaluations on one framed representative of each class. They coincide by equivariance on the remaining points.

The full-level class sum of that evaluation is \(1/|P_0|\)
times its principal-level class sum. For each full-level
orbit, with stabilizer \(S\), its principal-level classes
are the orbits of \(S\) on the finite level quotient
of size \(|P_0|\). Let \(X\) be this finite quotient. Orbit--stabilizer gives
\[
 \sum_{[x]\in S\backslash X}\frac1{|S_x|}
    =\frac1{|S|}\sum_{[x]\in S\backslash X}|Sx|
    =\frac{|P_0|}{|S|}.
\]
The Hermitian squared values are constant on this set
of classes because the level characters have modulus
one. The principal-level sum is \(\|VF\|_H^2\), so
the full-level evaluation has squared norm \(\|F\|^2\).
This gives the norm-nonincreasing evaluation required
in the terminal bound at every infinite place.
\label{cor:selected-lift-scalar-7}

Now use the quotient of this framed test space by the
kernel of its map \(q_0\) to the primitive rows. Its
metric is the norm of the orthogonal representative.
By \eqref{eq:selected-pairing} the terminal pairing
is constant on the fibres of \(q_0\), so it has the
same norm bound in the quotient metric. In the period
comparison the row isomorphism \(\iota\) from this
quotient to the spectral row multiplies its Hermitian
pairing with first argument of type \(p\) by
\[
             \kappa(v)\prod_{j=1}^m\lambda_{j,p}(v).
\]
This is the scalar identity for the same reference
row and word. The other argument is unrestricted.
Thus the type spaces lie in eigenspaces of the
positive self-adjoint comparison of the two metrics.
They span, and the least scalar is a lower bound for
the squared norm ratio on the whole quotient.

Passing to the spectral row dual divides the squared
operator bound by this least scalar. If \(H_j(v)\)
is the logarithmic squared kernel bound including
class sums, the resulting logarithmic squared bound
for a nonzero period map for the sequence \(B_u\) at \(v\) is
\[
 h_A(u)+\sum_{j=1}^m H_j(v)-\log\kappa(v)
                       -\min_p\sum_{j=1}^m\log\lambda_{j,p}(v).
\]
Here the two terminal kernel factors in the coupled
map are counted at their separate lift indices.
Average over the infinite places. The kernel bounds
of Section~\ref{subsec:selected-kernel-jets} and their
CM means give the prices \(h_B(j)\). Using
\eqref{eq:selected-hX} gives
\begin{equation}\label{eq:selected-reference-height}
 \operatorname{ht}(B_u)
    \le h_{\rm num}(u)-\bar\kappa,\qquad
 h_{\rm num}(u)=h_A(u)+\sum_{j=1}^m\bigl(h_B(j)-h_X(j)\bigr).
\end{equation}
The same numerical upper bound may be assigned to a zero period map for the sequence since its operator norm is zero.
The ranks give \(n_V\le N_m\), and \(L'=O(L_a+L)\)
by the selected type-prime and discriminant bounds.

Consequently the allowance in \(h_X(j)\) costs at
most \(O(fN_m^2(L_a+L))\) per lift, uniformly in the
word, place and type. Its hook main term is common
to all sequences on this family. Equation
\eqref{eq:selected-reference-height} uses the exact
numerical prices of the selected comparison after
increasing their one common absolute constant.
\label{cor:selected-lift-scalar-8}
\end{proof}


